\chapter{Funções contínuas} \label{lim:chapter}

%%%%%%%%%%%%%%%%%%%%%%%%%%%%%%%%%%%%%%%%%%%%%%%%%%%%%%%%%%%%%%%%%%%%%%%%%%%%%%

\section{Limites de funções}
\label{sec:limoffunc}

%mbxINTROSUBSECTION

\sectionnotes{2--3 aulas}

Antes de definirmos a continuidade de funções, consideremos uma noção de limite
um pouco mais geral que a de limite de uma sequência.
Dada uma função $f \colon S \to
\R$, queremos observar como $f(x)$ se comporta quando $x$ tende a certo ponto.

\subsection{Pontos de acumulação}

Primeiro,
retomemos um conceito que já vimos em um exercício.
Ao nos movermos dentro do conjunto $S$, só podemos nos aproximar de pontos
que tenham elementos de $S$ arbitrariamente próximos.

\begin{defn}
Seja $S \subset \R$ um conjunto. Um número $x \in \R$ é chamado
\emph{\myindex{ponto de acumulação}} de $S$
se, para todo $\epsilon > 0$, o conjunto $(x-\epsilon,x+\epsilon) \cap S
\setminus \{ x \}$ não é vazio.
\end{defn}

Isto é, $x$ é um ponto de acumulação de $S$ se há pontos de $S$
arbitrariamente próximos de $x$. Outra maneira de enunciar a definição é dizer
que $x$ é um ponto de acumulação de $S$ se, para todo $\epsilon > 0$, existe
$y \in S$ tal que $y \neq x$ e $\sabs{x - y} < \epsilon$.
Observe que um ponto de acumulação de $S$ não precisa pertencer a $S$.

Vejamos alguns exemplos.
\begin{enumerate}[(i)]
\item O conjunto
$\{ \nicefrac{1}{n} : n \in \N \}$ tem zero como seu único ponto de acumulação.
\item Os pontos de acumulação do intervalo aberto $(0,1)$ são
todos os pontos do intervalo fechado $[0,1]$.
\item
O conjunto dos pontos de acumulação de $\Q$ é toda a reta real $\R$.
\item
O conjunto dos pontos de acumulação de $[0,1) \cup \{ 2 \}$ é o intervalo $[0,1]$.
\item O conjunto $\N$ não tem pontos de acumulação em $\R$.
\end{enumerate}

\begin{prop}
Seja $S \subset \R$. Então $x \in \R$ é um ponto de acumulação de $S$
se, e somente se,
existe uma sequência convergente de números $\{ x_n \}_{n=1}^\infty$ tal que
$x_n \neq x$ e $x_n \in S$ para todo $n$, e $\lim\limits_{n\to\infty} x_n = x$.
\end{prop}

\begin{proof}
Suponhamos primeiro que $x$ seja um ponto de acumulação de $S$.
Para todo $n \in \N$, escolha $x_n$ como um ponto arbitrário de
$(x-\nicefrac{1}{n},x+\nicefrac{1}{n}) \cap S \setminus \{x\}$, que
não é vazio, pois $x$ é um ponto de acumulação de $S$.
Então
$x_n$ está a uma distância menor que $\nicefrac{1}{n}$ de $x$, isto é,
\begin{equation*}
\sabs{x-x_n} < \nicefrac{1}{n} .
\avoidbreak
\end{equation*}
Como $\{ \nicefrac{1}{n} \}_{n=1}^\infty$ converge para zero,
$\{ x_n \}_{n=1}^\infty$ converge para $x$.

Por outro lado, se começarmos com uma sequência de números
$\{ x_n \}_{n=1}^\infty$ em $S$
que converge para $x$ e tal que $x_n \neq x$ para todo $n$, então, para todo
$\epsilon > 0$, existe $M$ tal que, em particular, $\sabs{x_M - x} <
\epsilon$. Isto é, $x_M \in (x-\epsilon,x+\epsilon) \cap S \setminus \{x\}$.
\end{proof}

\subsection{Limites de funções}

Se uma função $f$ está definida em um conjunto $S$ e $c$ é um ponto de
acumulação de $S$, definimos então o limite de $f(x)$ quando $x$ se aproxima de $c$.
Para essa definição, é irrelevante se $f$ está definida em $c$ ou não.
Mesmo que a função esteja definida em $c$, o limite da função quando $x$ tende a
$c$ pode muito bem ser diferente de $f(c)$.

\begin{defn}
\index{limite!de uma função}%
Seja $f \colon S \to \R$ uma função e seja $c$ um ponto de acumulação de
$S \subset \R$.
Suponha que exista $L \in \R$ tal que, para todo $\epsilon > 0$,
exista $\delta > 0$ de modo que, sempre que $x \in S \setminus \{ c \}$
e $\sabs{x - c} < \delta$, tenhamos
\begin{equation*}
\babs{f(x) - L} < \epsilon .
\end{equation*}
Dizemos então que $f(x)$ \emph{converge}\index{convergente!função} para $L$ quando $x$ tende
a $c$, e escrevemos
\glsadd{not:limitasarrows}%
\begin{equation*}
f(x) \to L \quad\text{quando}\quad x \to c .
\end{equation*}
Dizemos que $L$ é um \emph{limite} de $f(x)$ quando $x$
tende a $c$ e, se $L$ é único (como de fato é), escrevemos
\glsadd{not:limfunc}%
\begin{equation*}
\lim_{x \to c} f(x) \coloneqq L .
\end{equation*}
Se não existe tal $L$, dizemos que o limite não existe ou
que $f$ \emph{\myindex{diverge}} em $c$.
\end{defn}

Mais uma vez, a notação e a linguagem usadas acima pressupõem que o limite $L$, se
existir, é único, o que precisa ser demonstrado. Observe que o fato de $c$
ser um ponto de acumulação é importante para demonstrar a unicidade.

\begin{prop}
Seja $c$ um ponto de acumulação de $S \subset \R$ e seja $f \colon S \to \R$
uma função tal que $f(x)$ converge quando $x$ tende a $c$. Então
o limite de $f(x)$ quando $x$ tende a $c$ é único.
\end{prop}

\begin{proof}
Sejam $L_1$ e $L_2$ dois números que satisfazem a definição.
Fixe $\epsilon > 0$ e encontre $\delta_1 > 0$ tal que
$\babs{f(x)-L_1} < \nicefrac{\epsilon}{2}$
para todo $x \in S \setminus \{c\}$ com $\sabs{x-c} < \delta_1$.
Encontre também $\delta_2 > 0$ tal que
$\babs{f(x)-L_2} < \nicefrac{\epsilon}{2}$
para todo $x \in S \setminus \{c\}$ com $\sabs{x-c} < \delta_2$.
Defina $\delta \coloneqq \min \{ \delta_1, \delta_2 \}$. Suponha que $x \in S$,
$\sabs{x-c} < \delta$ e $x \neq c$. Como $\delta > 0$ e $c$ é um ponto de acumulação,
existe tal $x$. Então
\begin{equation*}
\sabs{L_1 - L_2} =
\babs{L_1 - f(x) + f(x) - L_2} \leq
\babs{L_1 - f(x)} + \babs{f(x) - L_2} < \frac{\epsilon}{2} + \frac{\epsilon}{2}
= \epsilon.
\end{equation*}
Como $\sabs{L_1-L_2} < \epsilon$ para todo $\epsilon > 0$, temos
$L_1 = L_2$.
\end{proof}

\begin{example}
Considere $f \colon \R \to \R$ definida por $f(x) \coloneqq x^2$. Então,
para todo $c \in \R$,
\begin{equation*}
\lim_{x\to c} f(x) = \lim_{x\to c} x^2 = c^2 .
\end{equation*}

Demonstração: seja $c \in \R$ fixo e suponha que seja dado $\epsilon > 0$. Escreva
\begin{equation*}
\delta \coloneqq \min \left\{ 1 , \, \frac{\epsilon}{2\sabs{c}+1} \right\} .
\end{equation*}
Tome $x \neq c$ tal que $\sabs{x-c} < \delta$. Em particular,
$\sabs{x-c} < 1$. Pela desigualdade triangular reversa,
\begin{equation*}
\sabs{x}-\sabs{c} \leq \sabs{x-c} < 1 .
\end{equation*}
Somando $2\sabs{c}$ aos dois lados, obtemos
$\sabs{x} + \sabs{c} < 2\sabs{c} + 1$. Estimamos
\begin{equation*}
\begin{split}
\babs{f(x) - c^2} &= \sabs{x^2-c^2} \\
&= \babs{(x+c)(x-c)} \\
&= \sabs{x+c}\sabs{x-c} \\
&\leq \bigl(\sabs{x}+\sabs{c}\bigr)\sabs{x-c} \\
&< \bigl(2\sabs{c}+1\bigr)\sabs{x-c} \\
&< \bigl(2\sabs{c}+1\bigr)\frac{\epsilon}{2\sabs{c}+1} = \epsilon .
\end{split}
\end{equation*}
\end{example}

\begin{example}
Defina $f \colon [0,1) \to \R$ por
\begin{equation*}
f(x) \coloneqq
\begin{cases}
x & \text{se } x > 0 , \\
1 & \text{se } x = 0 .
\end{cases}
\end{equation*}
Então
$\lim\limits_{x\to 0} f(x) = 0$,
embora $f(0) = 1$. Veja \figureref{fig:limvaldiff}.
\begin{myfigureht}
\myincludegraphics{limvaldiff}{%
Gráfico de uma função que vale 1 quando a variável independente x é 0.
Quando a variável independente x é maior que 0, a função é igual a x.
Indica-se que a função não está definida em x=1 nem para valores de x
maiores que 1, nem para valores de x menores que 0.}
\caption{Função cujo limite em $0$ difere de seu valor.\label{fig:limvaldiff}}
\end{myfigureht}

Demonstração: seja dado $\epsilon > 0$. Defina $\delta \coloneqq \epsilon$.
Para $x \in [0,1)$, $x \neq 0$ e $\sabs{x-0} < \delta$, temos
\begin{equation*}
\babs{f(x) - 0} = \sabs{x} < \delta = \epsilon .
\end{equation*}
\end{example}

\subsection{Limites sequenciais} \label{subseq:sequentiallimits}

Relacionemos o limite definido acima com os limites de sequências.

\begin{lemma}\label{seqflimit:lemma}
Seja $S \subset \R$, seja $c$ um ponto de acumulação de $S$, seja $f \colon S \to
\R$ uma função e seja $L \in \R$.

Então
$f(x) \to L$ quando $x \to c$ se, e somente se, para toda sequência
$\{ x_n \}_{n=1}^\infty$
tal que $x_n \in S \setminus \{c\}$ para todo $n$
e $\lim_{n\to\infty} x_n = c$,
a sequência $\bigl\{ f(x_n) \bigr\}_{n=1}^\infty$ converge para $L$.
\end{lemma}

\begin{proof}
Suponhamos que
$f(x) \to L$ quando $x \to c$ e que $\{ x_n \}_{n=1}^\infty$ seja uma sequência
tal que
$x_n \in S \setminus \{c\}$ e
$\lim_{n\to\infty} x_n = c$.
Queremos mostrar que $\bigl\{ f(x_n) \bigr\}_{n=1}^\infty$ converge para $L$.
Seja dado $\epsilon > 0$. Encontre $\delta > 0$ tal que, se
$x \in S \setminus \{c\}$ e $\sabs{x-c} < \delta$, então
$\babs{f(x) - L} < \epsilon$. Como
$\{ x_n \}_{n=1}^\infty$ converge para $c$, encontre $M$ tal que, para $n \geq M$,
temos $\sabs{x_n - c} < \delta$. Portanto, para $n \geq M$,
\begin{equation*}
\babs{f(x_n) - L} < \epsilon .
\end{equation*}
Assim, $\bigl\{ f(x_n) \bigr\}_{n=1}^\infty$ converge para $L$.

Para a outra direção, usamos a demonstração por contraposição. Suponhamos
que não seja verdade que $f(x) \to L$ quando $x \to c$. A negação da
definição afirma que existe $\epsilon > 0$ tal que, para todo
$\delta > 0$, existe $x \in S \setminus \{c\}$ para o qual
$\sabs{x-c} < \delta$
e $\babs{f(x)-L} \geq \epsilon$.

Usemos $\nicefrac{1}{n}$ no lugar de $\delta$ na afirmação acima para
construir uma sequência $\{ x_n \}_{n=1}^\infty$. Existe
$\epsilon > 0$ tal que, para todo $n$,
existe um ponto $x_n \in S \setminus \{c\}$ para o qual
$\sabs{x_n-c} < \nicefrac{1}{n}$
e $\babs{f(x_n)-L} \geq \epsilon$.
A sequência $\{ x_n \}_{n=1}^\infty$ assim construída converge para $c$, mas
a sequência $\bigl\{ f(x_n) \bigr\}_{n=1}^\infty$ não converge para $L$.
E isso conclui a demonstração.
\end{proof}

É possível fortalecer a direção recíproca do lema simplesmente afirmando que
\myquote{$\bigl\{ f(x_n) \bigr\}_{n=1}^\infty$ converge,} sem exigir um limite específico.
Veja \exerciseref{exercise:seqflimitalt}.

\begin{example}
$\displaystyle \lim_{x \to 0} \sin( \nicefrac{1}{x} )$
não existe, mas
$\displaystyle \lim_{x \to 0} x\sin( \nicefrac{1}{x} ) = 0$.
Veja \figureref{figsin1x}.

\begin{myfigureht}
%left guy also used in 10.4
\myincludepdft{sin1x_xsin1x}{%
Os gráficos de funções em um intervalo ao redor da origem.
À esquerda, a função oscila entre menos 1 e 1
cada vez mais rapidamente à medida que x se aproxima de 0, a ponto de,
perto da origem, o gráfico parecer preto sólido.
À direita, a função também oscila cada vez mais rapidamente
à medida que x se aproxima de 0, mas agora a amplitude também diminui até 0,
de modo que a região que parece preta sólida se assemelha, na verdade, à
ponta de um triângulo.}
\caption{Gráficos de $\sin(\nicefrac{1}{x})$ e $x \sin(\nicefrac{1}{x})$.
Observe que o computador não consegue representar corretamente
$\sin(\nicefrac{1}{x})$ perto de zero, pois a oscilação é rápida demais.\label{figsin1x}}
\end{myfigureht}

Demonstração:
Comecemos com $\sin(\nicefrac{1}{x})$. Defina a sequência
por
$x_n \coloneqq \frac{1}{\pi n + \nicefrac{\pi}{2}}$. É fácil ver
que $\lim_{n\to\infty} x_n = 0$. Além disso,
\begin{equation*}
\sin ( \nicefrac{1}{x_n} )
=
\sin (\pi n + \nicefrac{\pi}{2})
= {(-1)}^n .
\end{equation*}
Portanto, a sequência $\bigl\{ \sin ( \nicefrac{1}{x_n} ) \bigr\}_{n=1}^\infty$ não converge.
Pelo
\lemmaref{seqflimit:lemma},
$\displaystyle \lim_{x \to 0} \sin( \nicefrac{1}{x} )$ não existe.

Agora consideremos $x\sin(\nicefrac{1}{x})$. Seja $\{ x_n \}_{n=1}^\infty$ uma sequência
tal que $x_n \neq 0$ para todo $n$ e $\lim_{n\to\infty} x_n = 0$.
Observe que $\babs{\sin(t)} \leq 1$ para todo $t \in \R$. Portanto,
\begin{equation*}
\babs{x_n\sin(\nicefrac{1}{x_n})-0}
=
\sabs{x_n}\babs{\sin(\nicefrac{1}{x_n})}
\leq
\sabs{x_n} .
\end{equation*}
Como $x_n$ tende a zero, $\sabs{x_n}$ também tende a zero e, portanto,
$\bigl\{ x_n\sin(\nicefrac{1}{x_n}) \bigr\}_{n=1}^\infty$ converge para zero. Pelo
\lemmaref{seqflimit:lemma},
$\displaystyle \lim_{x \to 0} x\sin( \nicefrac{1}{x} ) = 0$.
\end{example}

Tenha em mente a expressão \myquote{para toda sequência} no lema.
Por exemplo, tome $\sin(\nicefrac{1}{x})$ e a sequência dada por
$x_n \coloneqq \nicefrac{1}{\pi n}$.
Então $\bigl\{ \sin (\nicefrac{1}{x_n}) \bigr\}_{n=1}^\infty$
é a sequência constante zero e, portanto, converge para zero, mas o limite de
$\sin(\nicefrac{1}{x})$ quando $x \to 0$ não existe.

Usando o \lemmaref{seqflimit:lemma},
podemos começar a aplicar tudo o que sabemos sobre
limites de sequências aos limites de funções. Vejamos alguns exemplos importantes.

\begin{cor}
Seja $S \subset \R$ e seja $c$ um ponto de acumulação de $S$.
Suponha que $f \colon S \to
\R$ e $g \colon S \to \R$ sejam funções tais que
os limites de $f(x)$ e $g(x)$ quando $x$ tende a $c$ existam, e que
\begin{equation*}
f(x) \leq g(x) \qquad \text{para todo } x \in S \setminus \{ c \}.
\end{equation*}
Então
\begin{equation*}
\lim_{x\to c} f(x) \leq \lim_{x\to c} g(x) .
\end{equation*}
\end{cor}

\begin{proof}
Seja $\{ x_n \}_{n=1}^\infty$ uma sequência de números em $S \setminus \{ c \}$
que converge para $c$. Seja
\begin{equation*}
L_1 \coloneqq \lim_{x\to c} f(x), \qquad \text{e} \qquad L_2 \coloneqq \lim_{x\to c} g(x) .
\end{equation*}
O \lemmaref{seqflimit:lemma} afirma que $\bigl\{ f(x_n) \bigr\}_{n=1}^\infty$ converge para
$L_1$ e que $\bigl\{ g(x_n) \bigr\}_{n=1}^\infty$ converge para $L_2$. Também
temos $f(x_n) \leq g(x_n)$ para todo $n$.
Obtemos $L_1 \leq L_2$ usando o
\lemmaref{limandineq:lemma}.
\end{proof}

Aplicando funções constantes, obtemos o seguinte corolário. A demonstração
fica como exercício.

\begin{cor} \label{fconstineq:cor}
Seja $S \subset \R$ e seja $c$ um ponto de acumulação de $S$. Suponha que $f \colon S \to
\R$ seja uma função tal que o limite de $f(x)$ quando $x$ tende a $c$
exista.
Suponha que existam dois números reais $a$ e $b$ tais que
\begin{equation*}
a \leq f(x) \leq b \qquad \text{para todo } x \in S \setminus \{ c \}.
\end{equation*}
Então
\begin{equation*}
a \leq \lim_{x\to c} f(x) \leq b .
\end{equation*}
\end{cor}

Usando o \lemmaref{seqflimit:lemma} da mesma maneira que acima, obtemos também
os seguintes corolários, cujas demonstrações ficam como exercícios.

\begin{cor} \label{fsqueeze:cor}
Seja $S \subset \R$ e seja $c$ um ponto de acumulação de $S$.
Suponha que $f \colon S \to \R$,
$g \colon S \to \R$ e $h \colon S \to \R$ sejam funções tais que
\begin{equation*}
f(x) \leq g(x) \leq h(x) \qquad \text{para todo } x \in S \setminus \{ c \}.
\end{equation*}
Suponha que os limites de $f(x)$ e $h(x)$ quando $x$ tende a $c$ existam e que
\begin{equation*}
\lim_{x\to c} f(x) = \lim_{x\to c} h(x) .
\end{equation*}
Então o limite de $g(x)$ quando $x$ tende a $c$ existe e
\begin{equation*}
\lim_{x\to c} g(x) =
\lim_{x\to c} f(x) = \lim_{x\to c} h(x) .
\end{equation*}
\end{cor}

\begin{cor} \label{falg:cor}
Seja $S \subset \R$ e seja $c$ um ponto de acumulação de $S$.
Suponha que $f \colon S \to \R$ e
$g \colon S \to \R$ sejam funções tais que
os limites de $f(x)$ e $g(x)$ quando $x$ tende a $c$ existam.
Então
\begin{enumerate}[(i)]
\item
$\displaystyle
\lim_{x\to c} \bigl(f(x)+g(x)\bigr) = \left(\lim_{x\to c} f(x)\right) +
\left(\lim_{x\to c} g(x)\right) .
$
\item
$\displaystyle
\lim_{x\to c} \bigl(f(x)-g(x)\bigr) = \left(\lim_{x\to c} f(x)\right) -
\left(\lim_{x\to c} g(x)\right) .
$
\item
$\displaystyle
\lim_{x\to c} \bigl(f(x)g(x)\bigr) = \left(\lim_{x\to c} f(x)\right)
\left(\lim_{x\to c} g(x)\right) .
$
\item \label{falg:cor:iv} Se
$\displaystyle \lim_{x\to c} g(x) \neq 0$
e $g(x) \neq 0$ para todo $x \in S \setminus \{ c \}$, então
\begin{equation*}
\lim_{x\to c} \frac{f(x)}{g(x)} =
\frac{\lim_{x\to c} f(x)}{\lim_{x\to c} g(x)} .
\end{equation*}
\end{enumerate}
\end{cor}

\begin{cor} \label{fabs:cor}
Seja $S \subset \R$ e seja $c$ um ponto de acumulação de $S$.
Suponha que $f \colon S \to \R$ seja uma função tal que o limite de $f(x)$ quando $x$ tende a $c$
exista.
Então
\begin{equation*}
\lim_{x\to c} \babs{f(x)} =
\abs{\lim_{x\to c} f(x)}.
\end{equation*}
\end{cor}

\subsection{Limites de restrições e limites laterais}

Às vezes, trabalhamos com a função definida em um subconjunto.

\begin{defn}
Seja $f \colon S \to \R$ uma função e seja $A \subset S$. Defina a
função $f|_A \colon A \to \R$ por
\begin{equation*}
f|_A (x) \coloneqq f(x)  \qquad \text{para } x \in A.
\end{equation*}
Chamamos $f|_A$ de \emph{\myindex{restrição}} de $f$ a $A$.
\end{defn}

A função $f|_A$ é simplesmente a função $f$ considerada em um domínio menor.
A proposição a seguir é análoga a tomar uma cauda de uma sequência.
Ela afirma que o limite é \myquote{local}: o limite depende apenas dos pontos
arbitrariamente próximos de $c$.

\begin{prop} \label{prop:limrest}
Sejam $S \subset \R$, $c \in \R$ e
$f \colon S
\to \R$ uma função.
Suponha que
$A \subset S$ seja tal que, para algum $\alpha > 0$,
$(A \setminus \{ c \}) \cap (c-\alpha,c+\alpha) = (S \setminus \{ c \}) \cap (c-\alpha,c+\alpha)$.
\begin{enumerate}[(i)]
\item O ponto $c$ é um ponto de acumulação de $A$ se, e somente se, $c$ é um ponto de acumulação
de $S$.
\item Supondo que $c$ seja um ponto de acumulação de $S$, temos que
$f(x) \to L$ quando $x \to c$ se, e somente se,
$f|_A(x) \to L$ quando $x \to c$.
\end{enumerate}
\end{prop}

\begin{proof}
Primeiro, suponha que $c$ seja um ponto de acumulação de $A$.
Como $A \subset S$, se $( A \setminus \{ c\} ) \cap
(c-\epsilon,c+\epsilon)$ não é vazio para todo $\epsilon > 0$,
então $( S \setminus \{ c\} ) \cap
(c-\epsilon,c+\epsilon)$ também não é vazio para todo $\epsilon > 0$.
Portanto, $c$ é um ponto de acumulação de $S$.
Em segundo lugar, suponha que $c$ seja um ponto de acumulação
de $S$. Então, para $\epsilon > 0$ tal que $\epsilon < \alpha$,
temos $( A \setminus \{ c\} ) \cap (c-\epsilon,c+\epsilon) =
( S \setminus \{ c\} ) \cap (c-\epsilon,c+\epsilon)$, que é não vazio.
Isso vale para todo $\epsilon < \alpha$ e, portanto,
$( A \setminus \{ c\} ) \cap (c-\epsilon,c+\epsilon)$ deve ser não vazio para todo
$\epsilon > 0$. Assim, $c$ é um ponto de acumulação de $A$.

Agora suponha que $c$ seja um ponto de acumulação de $S$ e que $f(x) \to L$ quando $x \to c$.
Isto é, para todo $\epsilon > 0$, existe $\delta > 0$ tal que, se
$x \in S \setminus \{ c \}$
e $\sabs{x-c} < \delta$, então $\babs{f(x)-L} < \epsilon$. Como $A \subset S$,
se $x \in A \setminus \{ c \}$, então $x \in S \setminus \{ c \}$,
e, portanto, $f|_A(x) \to L$ quando $x \to c$.

Por fim, suponha que $f|_A(x) \to L$ quando $x \to c$ e seja dado $\epsilon > 0$.
Existe $\delta' > 0$ tal que, se $x \in A \setminus \{ c \}$
e $\sabs{x-c} < \delta'$, então $\babs{ f|_A(x)-L } < \epsilon$.
Defina $\delta \coloneqq \min \{ \delta', \alpha \}$.
Suponha agora que $x \in S \setminus \{ c \}$ e
$\sabs{x-c} < \delta$. Como $\sabs{x-c} < \alpha$, temos $x \in A \setminus \{ c \}$,
e, como $\sabs{x-c} < \delta'$,
obtemos $\babs{f(x)-L} = \babs{ f|_A(x)-L } < \epsilon$.
\end{proof}

A hipótese sobre $A$ na proposição é necessária. Para uma restrição
arbitrária, em geral obtemos uma implicação em apenas uma direção
(veja \exerciseref{exercise:restrictionlimitexercise}).
A notação usual para o limite é
\begin{equation*}
\lim_{\substack{x \to c\\x \in A}} f(x) \coloneqq \lim_{x \to c} f|_A(x) .
\end{equation*}
Um uso comum de restrições no estudo de limites, que não satisfaz
a hipótese da proposição, é o
\emph{\myindex{limite lateral}}%
\footnote{%
Há muitas notações para limites laterais.
Por exemplo,
$\lim\limits_{\substack{x \to c\\x < c}} f(x)$,
$\lim\limits_{x \uparrow c} f(x)$ ou
$\lim\limits_{x \nearrow c} f(x)$ para
$\lim\limits_{x \to c^-} f(x)$.}.

\begin{defn} \label{defn:onesidedlimits}
Seja $f \colon S \to \R$ uma função e seja $c \in \R$. Se
$c$ é um ponto de acumulação de
$S \cap (c,\infty)$
e existe o limite
da restrição de $f$ a $S \cap (c,\infty)$
quando $x \to c$, definimos
\glsadd{not:onesidedlim}
\begin{equation*}
\lim_{x \to c^+} f(x) \coloneqq \lim_{x\to c} f|_{S \cap (c,\infty)}(x) .
\end{equation*}
De modo semelhante, se $c$ é um ponto de acumulação de
$S \cap (-\infty,c)$ e existe o limite da restrição quando $x \to c$,
definimos
\begin{equation*}
\lim_{x \to c^-} f(x) \coloneqq \lim_{x\to c} f|_{S \cap (-\infty,c)}(x) .
\end{equation*}
\end{defn}

A \propref{prop:limrest} não se aplica a limites laterais.
Pode haver limites laterais sem que exista o limite em um ponto. Por
exemplo, defina $f \colon \R \to \R$ por $f(x) \coloneqq 1$ para $x < 0$ e
$f(x) \coloneqq 0$ para $x \geq 0$.
Deixamos ao leitor verificar que
\begin{equation*}
\lim_{x \to 0^-} f(x) = 1, \qquad
\lim_{x \to 0^+} f(x) = 0, \qquad
\lim_{x \to 0} f(x) \quad \text{não existe.}
\end{equation*}
Ainda assim, temos o seguinte resultado.

\begin{prop} \label{prop:onesidedlimits}
Seja $S \subset \R$ tal que $c$ seja um ponto de acumulação
tanto de $S \cap (-\infty,c)$ quanto de $S \cap (c,\infty)$, seja
$f \colon S \to \R$ uma função e seja $L \in \R$. Então $c$ é um ponto de acumulação de $S$ e
\begin{equation*}
\lim_{x \to c} f(x) = L
\qquad \text{se, e somente se,} \qquad
\lim_{x \to c^-} f(x) =
\lim_{x \to c^+} f(x) =
L .
\end{equation*}
\end{prop}

Isto é, o limite em $c$ existe se, e somente se, os dois limites laterais existem e são iguais. A
demonstração é uma aplicação direta da definição de limite e fica como exercício. O ponto central é que
$\bigl( S \cap (-\infty,c) \bigr) \cup \bigl( S \cap (c,\infty) \bigr)
= S \setminus \{ c \}$.



\subsection{Exercícios}

\begin{exercise}
Determine o limite (e, naturalmente, prove o resultado) ou prove que o limite não
existe:

\medskip

\noindent
\begin{tabular}{lllll}
a)
$\displaystyle
\lim_{x\to c} \sqrt{x}
$, para $c \geq 0$
& &
b)
$\displaystyle
\lim_{x\to c} x^2+x+1
$, for $c \in \R$
& &
c)
$\displaystyle
\lim_{x\to 0} x^2 \cos (\nicefrac{1}{x})
$
\\
d)
$\displaystyle
\lim_{x\to 0} \sin(\nicefrac{1}{x}) \cos (\nicefrac{1}{x})
$
& &
e)
$\displaystyle
\lim_{x\to 0} \sin(x) \cos (\nicefrac{1}{x})
$ & 
\end{tabular}
\end{exercise}

\begin{exercise}
Demonstre o \corref{fconstineq:cor}.
\end{exercise}

\begin{exercise}
Demonstre o \corref{fsqueeze:cor}.
\end{exercise}

\begin{exercise}
Demonstre o \corref{falg:cor}.
\end{exercise}

\begin{exercise}
Let $A \subset S$.  Mostre que if $c$ é um ponto de acumulação de $A$, then $c$
é um ponto de acumulação de $S$.  Observe a diferença em relação a
\propref{prop:limrest}.
\end{exercise}

\begin{exercise} \label{exercise:restrictionlimitexercise}
Seja $A \subset S$. Suponha que $c$ seja um ponto de acumulação de $A$ e
também de $S$.
Seja $f \colon S \to \R$ uma função. Mostre que, se
$f(x) \to L$ quando $x \to c$, então
$f|_A(x) \to L$ quando $x \to c$.
Observe a diferença em relação a
\propref{prop:limrest}.
\end{exercise}

\begin{exercise}
Encontre um exemplo de função $f \colon [-1,1] \to \R$ para a qual, tomando
$A\coloneqq [0,1]$, temos
$f|_A(x) \to 0$ quando $x \to 0$, mas o limite de $f(x)$ quando $x \to 0$
não existe. Explique por que não podemos aplicar
\propref{prop:limrest}.
\end{exercise}

\begin{exercise}
Encontre funções $f$ e $g$ tais que nem o limite de $f(x)$ nem o de $g(x)$ exista quando $x \to 0$, mas o limite de $f(x)+g(x)$ exista quando $x \to 0$.
\end{exercise}

\begin{exercise} \label{exercise:contlimitcomposition}
Seja $c_1$ um ponto de acumulação de $A \subset \R$ e $c_2$ um
ponto de acumulação de $B \subset \R$. Suponha que
$f \colon A \to B$ e $g \colon B \to \R$ sejam funções tais que
$f(x) \to c_2$ quando $x \to c_1$ e
$g(y) \to L$ quando $y \to c_2$. Se $c_2 \in B$, suponha também que $g(c_2) = L$.
Defina $h(x) \coloneqq g\bigl(f(x)\bigr)$ e mostre que
$h(x) \to L$ quando $x \to c_1$.
Dica: observe que $f(x)$ pode ser igual a $c_2$ para muitos $x \in A$
(veja também
\exerciseref{exercise:contlimitbadcomposition}).
\end{exercise}

\begin{exercise}
Suponha que $f \colon \R \to \R$ seja uma função tal que, para toda
sequência $\{x_n\}_{n=1}^\infty$ em $\R$, a sequência
$\bigl\{ f(x_n) \bigr\}_{n=1}^\infty$ converge. Prove que
$f$ é constante, isto é,
$f(x) = f(y)$ para todos $x,y \in \R$.
\end{exercise}

\begin{exercise} \label{exercise:seqflimitalt}
Demonstre a seguinte versão mais forte de uma das implicações do
\lemmaref{seqflimit:lemma}:
Seja $S \subset \R$, seja $c$ um ponto de acumulação de $S$ e seja $f \colon S \to
\R$ uma função.
Suponha que, para toda sequência $\{x_n\}_{n=1}^\infty$ em $S \setminus \{c\}$ tal que
$\lim_{n\to\infty} x_n = c$ a sequência $\bigl\{ f(x_n) \bigr\}_{n=1}^\infty$ seja convergente.
Mostre então que existe o limite de $f(x)$ quando $x \to c$.
\end{exercise}

\begin{exercise}
Demonstre a \propref{prop:onesidedlimits}.
\end{exercise}

\begin{exercise}
Suponha que $S \subset \R$ e que $c$ seja um ponto de acumulação de $S$.  Suponha que $f \colon
S \to \R$ seja limitada.  Mostre que existe uma sequência
$\{ x_n \}_{n=1}^\infty$
com $x_n \in S \setminus \{ c \}$ e $\lim_{n\to\infty} x_n = c$ tal que
$\bigl\{ f(x_n) \bigr\}_{n=1}^\infty$ converge.
\end{exercise}

\begin{exercise}[Desafio] \label{exercise:contlimitbadcomposition}
Mostre que a hipótese $g(c_2) = L$ em
\exerciseref{exercise:contlimitcomposition} é necessária. Isto é, encontre $f$
e $g$ tais que $f(x) \to c_2$ quando $x \to c_1$ e
$g(y) \to L$ quando $y \to c_2$, mas $g\bigl(f(x)\bigr)$ não tenda a $L$
quando $x \to c_1$.
\end{exercise}

\begin{exercise}
Mostre que the condition of being a cluster point is necessary to have a
reasonable definition of a limit.  That is, suppose $c$ is not a cluster
point of $S \subset \R$, and $f \colon S \to \R$ is a function.  Show that
every $L$ would satisfy the definition of limit at $c$ without the condition
on $c$ being a cluster point.
\end{exercise}

\begin{exercise}
\leavevmode
\begin{enumerate}[a)]
\item
Demonstre o \corref{fabs:cor}.
\item
Encontre um exemplo que mostre que a recíproca do
corolário não vale.
\end{enumerate}
\end{exercise}

%%%%%%%%%%%%%%%%%%%%%%%%%%%%%%%%%%%%%%%%%%%%%%%%%%%%%%%%%%%%%%%%%%%%%%%%%%%%%%




%%%%%%%%%%%%%%%%%%%%%%%%%%%%%%%%%%%%%%%%%%%%%%%%%%%%%%%%%%%%%%%%%%%%%%%%%%%%%%%

\section{Funções contínuas}
\label{sec:cont}

%mbxINTROSUBSECTION

\sectionnotes{2--2.5 aulas}

Um critério escolar para o conceito de continuidade é dizer que uma função é
contínua quando podemos desenhar seu gráfico sem tirar o lápis do papel.
Embora essa ideia intuitiva possa ser útil em situações simples, exigimos
rigor. Foram necessários três grandes matemáticos (Bolzano, Cauchy e,
finalmente, Weierstrass) para formular corretamente a definição a seguir; sua
forma final data apenas do fim do século XIX.

\subsection{Definição e propriedades básicas}

\begin{defn}
Sejam $S \subset \R$ e $c \in S$.
Dizemos que $f \colon S \to \R$
é \emph{contínua em $c$}\index{contínua em $c$}
se, para todo $\epsilon > 0$,
existe $\delta > 0$ tal que, sempre que $x \in S$ e $\sabs{x-c} <
\delta$, temos
$\babs{f(x)-f(c)} < \epsilon$.

%\medskip

Quando $f \colon S \to \R$ é contínua em todo $c \in S$, dizemos simplesmente
que $f$ é uma \emph{\myindex{função contínua}}\index{função!contínua}.
\end{defn}
\begin{myfigureht}
\myincludepdft{contigr}{%
Diagrama do gráfico de uma função y igual a f de x,
próximo de um ponto x igual a c. O gráfico é mostrado, assim como a
faixa a uma distância delta de c e a faixa a uma distância
epsilon de f de c. A caixa contida em ambas as faixas está sombreada;
nessa região, o gráfico fica dentro da faixa vertical.}
\caption{Quando $\sabs{x-c} < \delta$, o gráfico de $f(x)$ deve estar dentro da região cinza.\label{fig:contigr}}
\end{myfigureht}

Se $f$ é contínua para todo $c \in A$, dizemos que
$f$ é \emph{contínua em $A \subset S$}. Um exercício simples
(\exerciseref{exercise:restrictioncontinuous}) mostra que isso implica que
$f|_A$ é contínua, embora a recíproca não seja válida (como veremos no
\exampleref{example:removablediscont}).

Continuidade talvez seja a definição mais importante a compreender em análise,
e não é simples. Veja a \figureref{fig:contigr}.
Observe que $\delta$ depende não apenas
de $\epsilon$, mas também de $c$; não precisamos escolher
um único $\delta$ para todo $c \in S$.
Não é por acaso
que a definição de continuidade é semelhante à definição de limite de uma
função. A característica principal das funções contínuas
é que elas são precisamente as funções que se comportam bem com limites.

\begin{prop} \label{contbasic:prop}
Considere uma função $f \colon S \to \R$ definida em um conjunto $S \subset \R$
e seja $c \in S$.
Então:
%\begin{enumerate}[(i),itemsep=0.5\itemsep,parsep=0.5\parsep,topsep=0.5\topsep,partopsep=0.5\partopsep]
\begin{enumerate}[(i)]
\item \label{contbasic:prop:i}
Se $c$ não é ponto de acumulação de $S$, então $f$ é contínua em $c$.
\item \label{contbasic:prop:ii}
Se $c$ é ponto de acumulação de $S$, então $f$ é contínua em $c$
se, e somente se, existe o limite de $f(x)$ quando $x \to c$ e
\begin{equation*}
\lim_{x\to c} f(x) = f(c) .
\end{equation*}
\item \label{contbasic:prop:iii}
A função $f$ é contínua em $c$ se, e somente se, para toda sequência
$\{ x_n \}_{n=1}^\infty$
tal que $x_n \in S$ e $\lim\limits_{n\to\infty} x_n = c$, a sequência
$\bigl\{ f(x_n) \bigr\}_{n=1}^\infty$ converge
para $f(c)$.
\end{enumerate}
\end{prop}

\begin{proof}
\pagebreak[2]
Comecemos com \ref{contbasic:prop:i}. Suponha que $c$ não seja ponto de acumulação de
$S$. Então existe $\delta > 0$
tal que $S \cap (c-\delta,c+\delta) = \{ c \}$.
Para qualquer $\epsilon > 0$, basta escolher esse $\delta$.
O único $x \in S$ tal que $\sabs{x-c} < \delta$ é $x=c$. Portanto,
$\babs{f(x)-f(c)} = \babs{f(c)-f(c)} = 0 < \epsilon$.

Passemos a \ref{contbasic:prop:ii}.
Suponha que $c$ seja ponto de acumulação de $S$. Primeiro, suponhamos
que $\lim_{x\to c} f(x) = f(c)$. Então, para todo $\epsilon > 0$,
existe $\delta > 0$ tal que, se $x \in S \setminus \{ c \}$
e $\sabs{x-c} < \delta$, então $\babs{f(x)-f(c)} < \epsilon$.
Além disso, $\babs{f(c)-f(c)} = 0 < \epsilon$, de modo que a definição de continuidade em
$c$ é satisfeita. Por outro lado, suponha que $f$ seja contínua
em $c$. Para todo $\epsilon > 0$, existe $\delta > 0$
tal que, para $x \in S$ com $\sabs{x-c} < \delta$, temos
$\babs{f(x)-f(c)} < \epsilon$. Essa afirmação continua válida se
$x \in S \setminus \{ c \} \subset S$. Portanto, $\lim_{x\to c} f(x) =
f(c)$.

Para provar \ref{contbasic:prop:iii}, suponha primeiro que $f$ seja contínua em $c$.
Seja $\{ x_n \}_{n=1}^\infty$
uma sequência tal que $x_n \in S$ e $\lim_{n\to\infty} x_n = c$. Seja dado $\epsilon > 0$.
Encontre $\delta > 0$ tal que $\babs{f(x)-f(c)} < \epsilon$
para todo $x \in S$ com $\sabs{x-c} < \delta$. Encontre $M \in \N$
tal que, para $n \geq M$, tenhamos $\sabs{x_n-c} < \delta$. Então, para
$n \geq M$, temos $\babs{f(x_n)-f(c)} < \epsilon$,
e, portanto, $\bigl\{ f(x_n) \bigr\}_{n=1}^\infty$
converge para $f(c)$.

Provemos a outra implicação de \ref{contbasic:prop:iii} por contraposição.
Suponha que $f$ não seja contínua em $c$. Então existe $\epsilon > 0$
tal que, para todo $\delta > 0$, existe $x \in S$
tal que $\sabs{x-c} < \delta$ e $\babs{f(x)-f(c)} \geq \epsilon$.
Defina uma sequência $\{ x_n \}_{n=1}^\infty$ da seguinte maneira.
Escolha $x_n \in S$ tal que $\sabs{x_n-c} < \nicefrac{1}{n}$
e $\babs{f(x_n)-f(c)} \geq \epsilon$.
Então $\{ x_n \}_{n=1}^\infty$ é uma sequência em $S$ tal que
$\lim_{n\to\infty} x_n = c$ e $\babs{f(x_n)-f(c)} \geq \epsilon$ para todo $n \in \N$.
Assim, $\bigl\{ f(x_n) \bigr\}_{n=1}^\infty$
não converge para $f(c)$. Ela pode convergir ou não, mas certamente
não converge para $f(c)$.
\end{proof}

O último item da proposição é especialmente poderoso. Ele nos permite
aplicar rapidamente o que sabemos sobre limites de sequências a funções contínuas
e até provar que certas funções são contínuas.
Esse resultado também pode ser fortalecido (veja \exerciseref{exercise:contseqalt}).

\begin{example}
A função $f \colon (0,\infty) \to \R$
definida por $f(x) \coloneqq \nicefrac{1}{x}$ é contínua.

Demonstração: Fixe $c \in (0,\infty)$.
Seja $\{ x_n \}_{n=1}^\infty$ uma sequência em $(0,\infty)$ tal que
$\lim_{n\to\infty} x_n = c$. Então
\begin{equation*}
f(c) = \frac{1}{c}
=
\frac{1}{\lim_{n\to\infty} x_n}
=
\lim_{n \to \infty} \frac{1}{x_n}
=
\lim_{n \to \infty} f(x_n) .
\end{equation*}
Portanto, $f$ é contínua em $c$. Como $f$ é contínua em todo $c \in
(0,\infty)$, $f$ é contínua.
\end{example}

Já mostramos diretamente que $\lim_{x \to c} x^2 = c^2$. Portanto,
a função $x^2$ é contínua. O último item da \propref{contbasic:prop} e a continuidade das
operações algébricas em relação aos limites de sequências,
\propref{prop:contalg}, fornecem uma demonstração rápida de um resultado muito mais geral.

\begin{prop}
Seja $f \colon \R \to \R$ um \emph{\myindex{polinômio}}. Isto é,
\begin{equation*}
f(x) = a_d x^d + a_{d-1} x^{d-1} + \cdots + a_1 x + a_0 ,
\end{equation*}
para certas constantes $a_0, a_1, \ldots, a_d$.
Então $f$ é contínua.
\end{prop}

\begin{proof}
Fixe $c \in \R$.
Seja $\{ x_n \}_{n=1}^\infty$ uma sequência tal que
$\lim_{n\to\infty} x_n = c$. Então
\begin{equation*}
\begin{split}
f(c) &=
a_d c^d + a_{d-1} c^{d-1} + \cdots + a_1 c + a_0 
\\
&= 
a_d {\left(\lim_{n\to\infty} x_n\right)}^d
+ a_{d-1} {\left(\lim_{n\to\infty} x_n\right)}^{d-1}
+ \cdots
+ a_1 \left(\lim_{n\to\infty} x_n\right) + a_0 
\\
& =
\lim_{n \to \infty}
\left(
a_d x_n^d + a_{d-1} x_n^{d-1} + \cdots + a_1 x_n + a_0 
\right)
=
\lim_{n \to \infty}
f(x_n) .
\avoidbreak
\end{split}
\end{equation*}
Portanto, $f$ é contínua em $c$. Como $f$ é contínua em todo $c \in \R$,
$f$ é contínua.
\end{proof}

Por raciocínio semelhante, ou recorrendo ao \corref{falg:cor},
podemos provar a seguinte proposição. A demonstração fica como exercício.

\begin{prop} \label{contalg:prop}
Sejam $f \colon S \to \R$ e $g \colon S \to \R$ funções
contínuas em $c \in S$.
\begin{enumerate}[(i)]
\item A função $h \colon S \to \R$ definida por
$h(x) \coloneqq f(x)+g(x)$ é contínua em $c$.
\item A função $h \colon S \to \R$ definida por
$h(x) \coloneqq f(x)-g(x)$ é contínua em $c$.
\item A função $h \colon S \to \R$ definida por
$h(x) \coloneqq f(x)g(x)$ é contínua em $c$.
\item Se $g(x) \neq 0$ para todo $x \in S$, a função $h \colon S \to \R$
definida por $h(x) \coloneqq \frac{f(x)}{g(x)}$ é contínua em $c$.
\end{enumerate}
\end{prop}

\begin{example} \label{sincos:example}
As funções $\sin(x)$ e $\cos(x)$ são contínuas.
Nos cálculos a seguir, usamos as identidades trigonométricas de soma em produto.
Também usamos os fatos simples de que
$\babs{\sin(x)} \leq \sabs{x}$, $\babs{\cos(x)} \leq 1$,
and $\babs{\sin(x)} \leq 1$.
\begin{equation*}
\begin{split}
\babs{\sin(x)-\sin(c)} & =
\abs{
2 \sin \left( \frac{x-c}{2} \right) \cos \left( \frac{x+c}{2} \right)
}
\\
& =
2
\abs{ \sin \left( \frac{x-c}{2} \right) }
\abs{ \cos \left( \frac{x+c}{2} \right) }
\\
& \leq
2
\abs{ \sin \left( \frac{x-c}{2} \right) }
\\
& \leq
2
\abs{ \frac{x-c}{2} }
= \sabs{x-c}
\end{split}
\end{equation*}
\begin{equation*}
\begin{split}
\babs{\cos(x)-\cos(c)} & =
\abs{
-2 \sin \left( \frac{x-c}{2} \right) \sin \left( \frac{x+c}{2} \right)
}
\\
& =
2
\abs{ \sin \left( \frac{x-c}{2} \right) }
\abs{ \sin \left( \frac{x+c}{2} \right) }
\\
& \leq
2
\abs{ \sin \left( \frac{x-c}{2} \right) }
\\
& \leq
2
\abs{ \frac{x-c}{2} }
= \sabs{x-c}
\end{split}
\end{equation*}

A afirmação de que $\sin$ e $\cos$ são contínuas segue tomando uma
sequência arbitrária $\{ x_n \}_{n=1}^\infty$ convergente para $c$, ou aplicando diretamente a
definição de continuidade. Os detalhes ficam a cargo do
leitor.
\end{example}

\subsection{Composition of continuous functions}

Você provavelmente já percebeu que uma das ferramentas básicas para
construir funções complicadas a partir de funções simples é a composição.
Lembre-se de que, para duas funções $f$ e $g$,
a composição $f \circ g$ é definida por
$(f \circ g)(x) \coloneqq f\bigl(g(x)\bigr)$.
A composição de
funções contínuas também é
contínua.

\begin{prop} \label{prop:compositioncont}
Sejam $A, B \subset \R$ e sejam $f \colon B \to \R$ e $g \colon A \to B$
funções. Se $g$ é contínua em $c \in A$ e
$f$ é contínua em $g(c)$, então $f \circ g \colon A \to \R$ é contínua
em $c$.
\end{prop}

\begin{proof}
Seja $\{ x_n \}_{n=1}^\infty$ uma sequência em $A$ tal que
$\lim_{n\to\infty} x_n = c$.
Como $g$ é contínua em $c$, temos que $\bigl\{ g(x_n) \bigr\}_{n=1}^\infty$ converge para $g(c)$.
Como $f$ é contínua em $g(c)$, temos que $\bigl\{ f\bigl(g(x_n)\bigr)
\bigr\}_{n=1}^\infty$ converge
para $f\bigl(g(c)\bigr)$.
Portanto, $f \circ g$ é contínua em $c$.
\end{proof}

\begin{example}
Afirmação: \emph{${\bigl(\sin(\nicefrac{1}{x})\bigr)}^2$ é uma função contínua em $(0,\infty)$.}

Demonstração: A função $\nicefrac{1}{x}$ é contínua em
$(0,\infty)$ e $\sin(x)$ é contínua em $(0,\infty)$ (na verdade,
em $\R$, mas $(0,\infty)$ é a imagem de $\nicefrac{1}{x}$).
Logo, a composição $\sin(\nicefrac{1}{x})$ é contínua. Além disso,
$x^2$ é contínua no intervalo $[-1,1]$ (a imagem de $\sin$).
Assim, a composição
${\bigl(\sin(\nicefrac{1}{x})\bigr)}^2$ é contínua em $(0,\infty)$.
\end{example}

\subsection{Discontinuous functions}

Quando $f$ não é contínua em $c$, dizemos
que $f$ é \emph{\myindex{descontínua}} em $c$, ou que tem uma
\emph{\myindex{descontinuidade}} em~$c$.
A proposição a seguir fornece um critério útil e decorre imediatamente
do terceiro item da \propref{contbasic:prop}.

\begin{prop}
Seja $f \colon S \to \R$ uma função e seja $c \in S$. Suponha
que exista uma sequência $\{ x_n \}_{n=1}^\infty$, com $x_n \in S$ para todo $n$,
e $\lim_{n\to\infty} x_n = c$,
tal que $\bigl\{ f(x_n) \bigr\}_{n=1}^\infty$
não converge para $f(c)$.
Então $f$ é descontínua em $c$.
\end{prop}

Novamente,
dizer que
$\bigl\{ f(x_n) \bigr\}_{n=1}^\infty$ não converge para $f(c)$
significa que essa sequência
ou não converge,
ou converge para algo diferente de $f(c)$.

\begin{example} \label{example:jumpdiscont}
A função $f \colon \R \to \R$ definida por
\begin{equation*}
f(x) \coloneqq 
\begin{cases}
-1 & \text{se } x < 0, \\
1 & \text{se } x \geq 0
\end{cases}
\end{equation*}
não é contínua em 0.

Demonstração: Considere $\{ \nicefrac{-1}{n} \}_{n=1}^\infty$, que converge para 0. Então
$f(\nicefrac{-1}{n}) = -1$ para todo $n$,
e, portanto,
$\lim_{n\to\infty} f(\nicefrac{-1}{n}) = -1$, mas $f(0) = 1$. Assim, a função
não é contínua em 0. Veja a
\figureref{fig:jumpdiscont}.

\begin{myfigureht}
\myincludegraphics{jumpdiscont}{%
Gráfico de uma função que vale menos 1 para x negativo
e 1 para x não negativo. Os pontos menos 1,
menos um meio, menos um terço etc. estão marcados no
eixo x, e os valores correspondentes no
gráfico estão indicados por pontos.}
\caption{Descontinuidade de salto. Os valores de
$f(\nicefrac{-1}{n})$ e $f(0)$ estão marcados como pontos pretos.\label{fig:jumpdiscont}}
\end{myfigureht}

Observe que $f(\nicefrac{1}{n}) = 1$ para todo $n \in \N$.
Logo, $\lim_{n\to\infty} f(\nicefrac{1}{n}) = f(0) = 1$. Portanto,
$\bigl\{ f(x_n) \bigr\}_{n=1}^\infty$ pode convergir para $f(0)$
para alguma sequência específica
$\{ x_n \}_{n=1}^\infty$ que tende a 0,
embora a função seja descontínua em 0.

Por fim, considere $f\Bigl(\frac{{(-1)}^n}{n}\Bigr) = {(-1)}^n$.
Essa sequência diverge.
\end{example}

\begin{example}
Como exemplo extremo, considere a
\emph{\myindex{função de Dirichlet}}\footnote{Nomeada em homenagem ao matemático alemão
\href{https://en.wikipedia.org/wiki/Peter_Gustav_Lejeune_Dirichlet}{Johann Peter Gustav Lejeune
Dirichlet}
(1805--1859).}.
\begin{equation*}
f(x) \coloneqq
\begin{cases}
1 & \text{se } x \text{ é racional,} \\
0 & \text{se } x \text{ é irracional.}
\end{cases}
\avoidbreak
\end{equation*}
A função $f$ é descontínua em todo $c \in \R$.

Demonstração:
Se $c$ é racional, tome uma sequência $\{ x_n \}_{n=1}^\infty$
de números irracionais em $(0,1)$ tal que $\lim_{n\to\infty} x_n = c$ (por que isso é possível?). Então $f(x_n) = 0$
e, portanto, $\lim_{n\to\infty} f(x_n) = 0$, mas $f(c) = 1$.
Se $c$ é irracional, tome uma sequência
$\{ x_n \}_{n=1}^\infty$
de números racionais em $(0,1)$
que converge para $c$ (por que isso é possível?). Então $\lim_{n\to\infty} f(x_n) = 1$, mas $f(c) = 0$.
\end{example}

Vamos testar os limites de nossa intuição. Pode
existir uma função contínua em todos os números irracionais, mas
descontínua em todos os racionais? Há racionais
arbitrariamente próximos de qualquer número irracional. Talvez surpreendentemente, a
resposta seja sim: tal função existe. O exemplo a seguir é chamado de
\emph{\myindex{função de Thomae}}\footnote{Nomeada em homenagem ao matemático alemão
\href{https://en.wikipedia.org/wiki/Carl_Johannes_Thomae}{Carl Johannes Thomae}
(1840--1921).} ou
\emph{\myindex{função pipoca}}.

\begin{example} \label{popcornfunction:example}
Defina $f \colon (0,1) \to \R$ por
\begin{equation*}
f(x) \coloneqq 
\begin{cases}
\nicefrac{1}{k} & \text{se } x=\nicefrac{m}{k}, \text{ where } m,k \in \N
\text{ and have no common divisors (lowest terms),} \\
0 & \text{se } x \text{ é irracional.}
\end{cases}
\end{equation*}
Veja o gráfico de $f$
na \figureref{popcornfig}.
Afirmamos que
$f$ é contínua em todo irracional $c$ e
descontínua em todo racional $c$.
\begin{myfigureht}
\myincludegraphics{popcornfig}{%
Gráfico de uma função formada por pontos.
Esses pontos parecem estar dispostos em
vários triângulos, mas é claro que a função
é descontínua em todos esses pontos,
e há cada vez mais pontos desse tipo
quanto mais nos aproximamos do eixo x.}
\caption{Gráfico da \myquote{função pipoca.}\label{popcornfig}}
\end{myfigureht}

Demonstração:
Seja $c = \nicefrac{m}{k}$ racional e escrito em termos irredutíveis. Tome uma sequência de
números irracionais $\{ x_n \}_{n=1}^\infty$ tal que $\lim_{n\to\infty} x_n = c$. Então
$\lim_{n\to\infty} f(x_n) = \lim_{n\to\infty} 0 = 0$, but $f(c) = \nicefrac{1}{k} \neq 0$.  Portanto, $f$
é descontínua em $c$.

Agora suponha que $c$ seja irracional, de modo que $f(c) = 0$. Tome uma sequência
$\{ x_n \}_{n=1}^\infty$ em $(0,1)$ tal que $\lim_{n\to\infty} x_n = c$.
Dado $\epsilon > 0$, encontre $K \in \N$ tal
que $\nicefrac{1}{K} < \epsilon$
pela \hyperref[thm:arch:i]{propriedade arquimediana}.
Se $\nicefrac{m}{k} \in (0,1)$ e $m,k \in \N$, então $0 < m < k$.
Assim, há apenas um número finito de racionais em $(0,1)$
cujo denominador $k$, na forma irredutível, é menor que $K$.
Como $\lim_{n\to\infty} x_n = c$, todo número diferente de $c$ pode aparecer no máximo
um número finito de vezes em $\{ x_n \}_{n=1}^\infty$.
Consequentemente,
existe $M$ tal que, para $n \geq M$, todos os termos racionais $x_n$
têm denominador maior ou igual a $K$. Portanto, para $n \geq M$,
\begin{equation*}
\babs{f(x_n) - 0} = f(x_n) \leq \nicefrac{1}{K} < \epsilon .
\end{equation*}
Portanto, $f$ é contínua em todo $c$ irracional.
\end{example}

Terminemos com um exemplo mais simples.

\begin{example} \label{example:removablediscont}
Defina
$g \colon \R \to \R$ por $g(x) \coloneqq 0$ se $x \neq 0$ e
$g(0) \coloneqq 1$. Então $g$ não é contínua em zero, mas é contínua em todos os outros pontos (por quê?).
O ponto $x=0$ é chamado de \emph{\myindex{descontinuidade removível}}. Isso
porque, se alterássemos a definição de $g$, estipulando que
$g(0)$ fosse $0$, obteríamos uma função contínua. Por outro lado,
seja $f$ a função do \exampleref{example:jumpdiscont}.
Então $f$ não tem uma
descontinuidade removível em $0$. Não importa como definamos $f(0)$, a função
ainda não será contínua. A diferença é que
$\lim_{x\to 0} g(x)$ existe, enquanto
$\lim_{x\to 0} f(x)$ não existe.

Continuemos com este exemplo para mostrar outro fenômeno.
Seja $A \coloneqq \{ 0 \}$; então $g|_A$ é contínua (por quê?), enquanto $g$ não é contínua em $A$.
De modo semelhante, se $B \coloneqq \R \setminus \{0 \}$, então $g|_B$ também é contínua,
e, de fato, $g$ é contínua em $B$.
\end{example}

\subsection{Exercises}

\begin{exercise}
Encontre um exemplo de função descontínua $f \colon [0,1] \to \R$
para a qual a conclusão do teorema do valor intermediário não vale.
\end{exercise}

\begin{exercise}
Encontre um exemplo de função descontínua \emph{limitada}
$f \colon [0,1] \to \R$ que não tenha mínimo absoluto nem máximo absoluto.
\end{exercise}

\begin{exercise}
Seja $f \colon (0,1) \to \R$ uma função contínua tal que
$\displaystyle \lim_{x\to 0} f(x) =
\displaystyle \lim_{x\to 1} f(x) = 0$. Mostre que
$f$ atinge um mínimo absoluto ou um máximo absoluto em $(0,1)$
(mas talvez não ambos).
\end{exercise}

\begin{exercise} \label{exercise:meanvaluepropsin1x}
Seja
\begin{equation*}
f(x) \coloneqq
\begin{cases}
\sin(\nicefrac{1}{x}) & \text{se } x \neq 0, \\
0 & \text{se } x=0.
\end{cases}
\end{equation*}
Mostre que $f$ tem a propriedade do valor intermediário.
Isto é, sempre que $a < b$, se existe $y$ tal que $f(a) < y < f(b)$
ou $f(a) > y > f(b)$, então
existe $c \in (a,b)$ tal que $f(c) = y$.
\end{exercise}

\begin{exercise} \label{exercise:odddegnegativeK}
Suponha que $g(x)$ seja um polinômio mônico de grau ímpar $d$, isto é,
\begin{equation*}
g(x) = x^d + b_{d-1} x^{d-1} + \cdots + b_1 x + b_0 ,
\end{equation*}
para certos números reais $b_{0}, b_1, \ldots, b_{d-1}$. Mostre que existe
$K \in \N$ tal que $g(-K) < 0$. Dica: certifique-se de usar o fato de que
$d$ é ímpar. Você precisará usar que ${(-n)}^d = -(n^d)$.
\end{exercise}

\begin{exercise}
Suponha que $g(x)$ seja um polinômio mônico de grau par positivo $d$, isto é,
\begin{equation*}
g(x) = x^d + b_{d-1} x^{d-1} + \cdots + b_1 x + b_0 ,
\end{equation*}
para certos números reais $b_{0}, b_1, \ldots, b_{d-1}$. Suponha
que $g(0) < 0$. Mostre que $g$ tem pelo menos duas raízes reais distintas.
\end{exercise}

\begin{exercise}
Demonstre o \corref{cor:imageofinterval}:
suponha que $f \colon [a,b] \to \R$ seja uma função contínua.
Demonstre que a imagem direta $f\bigl([a,b]\bigr)$ é um intervalo fechado
e limitado ou um único número.
\end{exercise}

\begin{exercise}
Suponha que $f \colon \R \to \R$ seja contínua e periódica, com período
$P > 0$. Isto é, $f(x+P) = f(x)$ para todo $x \in \R$. Mostre que $f$
atinge um mínimo absoluto e um máximo absoluto.
\end{exercise}

\begin{exercise}[Desafiador]
Suponha que $f(x)$ seja um polinômio limitado; em outras palavras,
existe $M$ tal que $\babs{f(x)} \leq M$
para todo $x \in \R$. Prove que $f$ deve ser constante.
\end{exercise}

\begin{exercise}
Suponha que $f \colon [0,1] \to [0,1]$ seja contínua. Mostre que $f$
tem um ponto fixo; em outras palavras, mostre que existe $x \in [0,1]$ tal que
$f(x) = x$.
\end{exercise}

\begin{exercise}
Encontre um exemplo de função contínua e limitada $f \colon \R \to \R$ que
não atinja mínimo absoluto nem máximo absoluto em $\R$.
\end{exercise}

\begin{exercise}
Suponha que $f \colon \R \to \R$ seja contínua e satisfaça
$x \leq f(x) \leq x+1$ para todo $x \in \R$. Determine $f(\R)$.
\end{exercise}

\begin{exercise}
Verdadeiro ou falso? Demonstre ou encontre um contraexemplo.
Se $f \colon \R \to \R$ é uma função contínua tal que $f|_{\Z}$ é limitada,
então $f$ é limitada.
\end{exercise}

\begin{exercise}
Suponha que $f \colon [0,1] \to (0,1)$ seja uma bijeção. Prove que $f$
não é contínua.
\end{exercise}

\begin{exercise}
Suponha que $f \colon \R \to \R$ seja contínua.
\begin{enumerate}[a)]
\item
Prove que, se existe $c$ tal que $f(c)f(-c) < 0$,
então existe $d \in \R$ tal que $f(d) = 0$.
\item
Encontre uma função contínua $f$ tal que
$f(\R) = \R$, mas $f(x)f(-x) \geq 0$ para todo $x \in \R$.
\end{enumerate}
\end{exercise}

\begin{exercise}
Suponha que $g(x)$ seja um polinômio mônico de grau par $d$, isto é,
\begin{equation*}
g(x) = x^d + b_{d-1} x^{d-1} + \cdots + b_1 x + b_0 ,
\end{equation*}
para certos números reais $b_{0}, b_1, \ldots, b_{d-1}$.
Mostre que $g$ atinge um mínimo absoluto em $\R$.
\end{exercise}

\begin{exercise}
Suponha que $f(x)$ seja um polinômio de grau $d$ e que
$f(\R) = \R$. Mostre que $d$ é ímpar.
\end{exercise}

%%%%%%%%%%%%%%%%%%%%%%%%%%%%%%%%%%%%%%%%%%%%%%%%%%%%%%%%%%%%%%%%%%%%%%%%%%%%%%

\sectionnewpage
\sectionnewpage
\section{Teoremas do valor extremo e do valor intermediário}
\label{sec:minmaxint}

%mbxINTROSUBSECTION

\sectionnotes{1,5 aulas}

Funções contínuas em intervalos fechados e limitados
têm um comportamento bastante regular.

\subsection{Teorema do máximo-mínimo ou do valor extremo}

Lembre-se de que $f \colon [a,b] \to \R$ é
\emph{limitada}\index{função limitada}\index{função!limitada}
se existe $B \in \R$ tal que
$\babs{f(x)} \leq B$ para todo $x \in [a,b]$.
Para uma função contínua em um intervalo fechado e limitado, temos o lema a seguir.

\begin{lemma}
Uma função contínua $f \colon [a,b] \to \R$ é limitada.
\end{lemma}

\begin{proof}
Provaremos a afirmação por contraposição. Suponha que $f$ não seja limitada.
Então, para cada
$n \in \N$, existe $x_n \in [a,b]$ tal que
\begin{equation*}
\babs{f(x_n)} \geq n .
\end{equation*}
A sequência $\{ x_n \}_{n=1}^\infty$ é limitada, pois $a \leq x_n \leq b$.
Pelo \hyperref[thm:bwseq]{teorema de Bolzano--Weierstrass},
existe uma subsequência convergente $\{ x_{n_i} \}_{i=1}^\infty$.
Seja $x \coloneqq \lim_{i\to\infty} x_{n_i}$.
Como $a \leq x_{n_i} \leq b$ para todo $i$, temos $a \leq x \leq b$.
A sequência $\bigl\{ f(x_{n_i}) \bigr\}_{i=1}^\infty$ não é limitada,
pois
$\babs{f(x_{n_i})} \geq n_i \geq i$.
Assim, $f$ não é contínua em $x$, pois
\begin{equation*}
f(x)
=
f\Bigl( \lim_{i\to\infty} x_{n_i} \Bigr) ,
\qquad \text{mas} \qquad
\lim_{i\to\infty} f(x_{n_i}) \enspace \text{não existe.} \qedhere
\end{equation*}
\end{proof}

Observe um ponto importante da demonstração.
O fato de $[a,b]$ ser limitado nos permite usar Bolzano--Weierstrass,
enquanto o fato de ser fechado garante que o limite continua em $[a,b]$.
A técnica é comum: encontrar uma sequência com certa propriedade
e então usar Bolzano--Weierstrass para obter uma subsequência que também convirja.

Relembre do cálculo que $f \colon S \to \R$ atinge um
\emph{\myindex{mínimo absoluto}}\index{mínimo!absoluto}%
\index{atinge mínimo absoluto}
em $c \in S$ se
\begin{equation*}
f(x) \geq f(c) \qquad \text{para todo } x \in S.
\end{equation*}
Por outro lado, $f$ atinge um
\emph{\myindex{máximo absoluto}}\index{máximo!absoluto}%
\index{atinge máximo absoluto}
em $c \in S$ se
\begin{equation*}
f(x) \leq f(c) \qquad \text{para todo } x \in S.
\end{equation*}
Se tal $c \in S$ existe, dizemos que
$f$ \emph{atinge um mínimo absoluto (respectivamente, máximo absoluto) em
$S$}, e chamamos
$f(c)$ de \emph{mínimo absoluto (respectivamente, máximo absoluto)}.
%See \figureref{fig:minmax}.
\begin{myfigureht}
\myincludegraphics{minmax}{%
Gráfico de uma função no intervalo de a até b.
O gráfico está entre duas retas horizontais,
uma indicada como o máximo absoluto de f, que é igual a
f de c, e
a outra indicada como o mínimo absoluto de f, que é igual a
f de d.}
\caption{$f \colon [a,b] \to \R$ atinge um máximo absoluto $f(c)$ em
$c$ e um mínimo absoluto $f(d)$ em $d$.\label{fig:minmax}}
\end{myfigureht}

Se $S$ é um intervalo fechado
e limitado, então uma função contínua $f$
não é apenas limitada: ela também atinge um mínimo absoluto e um máximo absoluto
em $S$.

\begin{thm}[Teorema do máximo-mínimo / Teorema do valor extremo]
\index{teorema do mínimo-máximo}%
\index{teorema do máximo-mínimo}%
\index{teorema do valor extremo}%
Uma função contínua $f \colon [a,b] \to \R$
atinge um mínimo absoluto e um máximo absoluto em $[a,b]$.
\end{thm}

Mais uma vez, ressaltamos que é importante que o domínio de $f$ seja um intervalo fechado e limitado $[a,b]$.

\begin{proof}
O lema afirma que $f$ é limitada; portanto,
o conjunto $f\bigl([a,b]\bigr) = \bigl\{ f(x) : x \in [a,b] \bigr\}$ tem supremo e ínfimo.
Existem sequências
no conjunto $f\bigl([a,b]\bigr)$ que se aproximam de seu supremo e de seu ínfimo.
Isto é, há sequências
$\bigl\{ f(x_n) \bigr\}_{n=1}^\infty$ e $\bigl\{ f(y_n)
\bigr\}_{n=1}^\infty$, em que $x_n$ e $y_n$ pertencem a $[a,b]$,
tais que
\begin{equation*}
\lim_{n\to\infty} f(x_n) = \inf f\bigl([a,b]\bigr) \qquad \text{e} \qquad
\lim_{n\to\infty} f(y_n) = \sup f\bigl([a,b]\bigr).
\end{equation*}
Ainda não terminamos: precisamos encontrar onde estão os mínimos e máximos.
O problema é que as sequências $\{ x_n \}_{n=1}^\infty$ e
$\{ y_n \}_{n=1}^\infty$ não precisam convergir.
Sabemos que $\{ x_n \}_{n=1}^\infty$ e $\{ y_n \}_{n=1}^\infty$ são limitadas
(seus termos pertencem ao intervalo limitado $[a,b]$).
Apliquemos o
\hyperref[thm:bwseq]{teorema de Bolzano--Weierstrass}
para encontrar
subsequências convergentes
$\{ x_{n_i} \}_{i=1}^\infty$ e
$\{ y_{m_i} \}_{i=1}^\infty$. Seja
\begin{equation*}
x \coloneqq \lim_{i\to\infty} x_{n_i}
\qquad \text{e} \qquad
y \coloneqq \lim_{i\to\infty} y_{m_i}.
\end{equation*}
Como $a \leq x_{n_i} \leq b$ para todo $i$, temos $a \leq x \leq b$.
De modo semelhante, $a \leq y \leq b$. Portanto, $x$ e $y$ pertencem a $[a,b]$.
Para sequências convergentes,
o limite de uma subsequência é igual ao limite da
sequência. Além disso, podemos passar o limite pela função contínua $f$. Assim,
\begin{equation*}
\inf f\bigl([a,b]\bigr) = \lim_{n\to\infty} f(x_n)
= \lim_{i\to\infty} f(x_{n_i}) = 
f \Bigl( \lim_{i\to\infty} x_{n_i} \Bigr) = f(x) .
\end{equation*}
De modo semelhante,
\begin{equation*}
\sup f\bigl([a,b]\bigr) = \lim_{n\to\infty} f(y_n)
= \lim_{i\to\infty} f(y_{m_i}) = 
f \Bigl( \lim_{i\to\infty} y_{m_i} \Bigr) = f(y) .
\end{equation*}
Portanto, $f$ atinge um mínimo absoluto em $x$ e
um máximo absoluto em $y$.
\end{proof}

\begin{example}
A função $f(x) \coloneqq x^2+1$ definida no intervalo $[-1,2]$ atinge um mínimo
em $x=0$, onde $f(0) = 1$. Atinge um máximo em $x=2$, onde $f(2) = 5$.
Observe que o domínio é relevante. Se, em vez disso, tomássemos o domínio
como $[-10,10]$, então $f$ não teria mais um máximo em $x=2$.
O máximo seria atingido em $x=10$ e também em $x=-10$.
\end{example}

Mostraremos com exemplos que as diferentes hipóteses do teorema são
realmente necessárias.

\begin{example}
A função $f \colon \R \to \R$ definida por $f(x) \coloneqq x$
não atinge mínimo nem máximo. Portanto, é importante
considerarmos um intervalo limitado.
\end{example}

\begin{example}
A função $f \colon (0,1) \to \R$ definida por $f(x) \coloneqq \nicefrac{1}{x}$
não atinge mínimo nem máximo.
Ela é contínua, mas $(0,1)$ não é fechado.
Os valores da função são
ilimitados quando nos aproximamos de $0$. Além disso, quando nos aproximamos de $x=1$, os valores da
função se aproximam de $1$, mas $f(x) > 1$ para todo $x \in (0,1)$. Não existe
$x \in (0,1)$ tal que $f(x) = 1$. Portanto, é importante
considerarmos um intervalo fechado.
\end{example}

\begin{example}
A continuidade é importante.
Defina $f \colon [0,1] \to \R$ por
$f(x) \coloneqq \nicefrac{1}{x}$ para $x > 0$ e seja $f(0) \coloneqq 0$.
A função não atinge um máximo. O domínio $[0,1]$ é fechado e
limitado, mas o problema é que
a função não é contínua em 0.
\end{example}

\subsection{Bolzano's intermediate value theorem}

O teorema do valor intermediário de Bolzano é um dos pilares da análise.
Às vezes, ele é chamado apenas de teorema do valor intermediário ou,
simplesmente, de teorema de Bolzano. Para demonstrar o teorema de Bolzano,
provaremos o lema mais simples a seguir.

\begin{lemma} \label{IVT:lemma}
Seja $f \colon [a,b] \to \R$ uma função contínua.
Suponha que $f(a) < 0$ e $f(b) > 0$.
Então existe um número $c \in (a,b)$ tal que $f(c) = 0$.
\end{lemma}

\begin{proof}
Definimos indutivamente duas sequências $\{ a_n \}_{n=1}^\infty$
e $\{ b_n \}_{n=1}^\infty$:
\begin{enumerate}[(i)]
\item Seja $a_1 \coloneqq a$ e $b_1 \coloneqq b$.
\item Se $f\left(\frac{a_n+b_n}{2}\right) \geq 0$, seja
$a_{n+1} \coloneqq a_n$ e
$b_{n+1} \coloneqq \frac{a_n+b_n}{2}$.
\item Se $f\left(\frac{a_n+b_n}{2}\right) < 0$, seja
$a_{n+1} \coloneqq \frac{a_n+b_n}{2}$ e
$b_{n+1} \coloneqq b_n$.
\end{enumerate}
\begin{myfigureht}
\myincludegraphics{bisect}{%
Gráfico de uma função que cruza o eixo x em c, crescendo.
O intervalo de a_1 a b_1 está assinalado, e f é negativa
em a_1 e positiva em b_1.
O intervalo seguinte tem a_2 igual a a_1, e b_2 está no meio
do intervalo anterior. Novamente, a função é negativa em a_2 e
positiva em b_2. Prosseguimos com os intervalos nos quais f é
negativa à esquerda e positiva à direita.
O intervalo seguinte vai de a_3, que é igual a a_2, até b_3,
que está no meio do intervalo anterior.
No intervalo seguinte, a_4 está no meio e b_4 é igual a b_3.
Depois, a_5 volta a estar no meio e b_5 é igual a b_4.
O número c está dentro de todos esses intervalos.}
\caption{Encontrando raízes (método da bisseção).\label{bisectfig}}
\end{myfigureht}
A \figureref{bisectfig} mostra um exemplo dos cinco primeiros passos.
Se $a_n < b_n$, então
$a_n < \frac{a_n+b_n}{2} < b_n$. Portanto,
$a_{n+1} < b_{n+1}$.
Como $a_1 = a < b = b_1$,
a \hyperref[induction:thm]{indução} implica que
$a_n < b_n$ para todo $n$.
Além disso, $a_n \leq a_{n+1}$ e
$b_n \geq b_{n+1}$ para todo $n$; isto é, as sequências são monótonas.
Como $a_n < b_n \leq b_1 = b$ e
$b_n > a_n \geq a_1 = a$ para todo $n$,
as sequências também são limitadas. Portanto, elas convergem.
Sejam $c \coloneqq \lim_{n\to\infty} a_n$ e
$d \coloneqq \lim_{n\to\infty} b_n$,
de modo que também $a \leq c \leq d \leq b$. Precisamos mostrar que $c=d$.
Observe que
\begin{equation*}
b_{n+1} - a_{n+1} = \frac{b_n-a_n}{2}.
\end{equation*}
Pela \hyperref[induction:thm]{indução},
\begin{equation*}
b_n - a_n = \frac{b_1-a_1}{2^{n-1}} = 2^{1-n} (b-a) .
\end{equation*}
Como $2^{1-n}(b-a)$ converge para zero, tomamos o limite quando $n$ tende
ao infinito e obtemos
\begin{equation*}
d-c = \lim_{n\to\infty} (b_n - a_n) =
\lim_{n\to\infty} 2^{1-n} (b-a) = 0.
\avoidbreak
\end{equation*}
Em outras palavras, $c=d$.

Por construção, para todo $n$,
\begin{equation*}
f(a_n) < 0
\qquad \text{e} \qquad
f(b_n) \geq 0 .
\end{equation*}
Como
$\lim_{n\to\infty} a_n = \lim_{n\to\infty} b_n = c$
e $f$ é contínua em $c$, podemos tomar limites nessas desigualdades:
\begin{equation*}
f(c) = \lim_{n\to\infty} f(a_n) \leq 0
\qquad \text{e} \qquad
f(c) = \lim_{n\to\infty} f(b_n) \geq 0 .
\end{equation*}
Como $f(c) \geq 0$ e $f(c) \leq 0$, concluímos que $f(c) = 0$.
Logo, também temos $c \neq a$ e $c \neq b$, portanto $a < c < b$.
\end{proof}

\begin{thm}[Teorema do valor intermediário de Bolzano] \label{IVT:thm}
\index{teorema de Bolzano}
\index{teorema do valor intermediário de Bolzano}
\index{teorema do valor intermediário}
Seja $f \colon [a,b] \to \R$ uma função contínua.
Suponha que $y \in \R$ seja tal que $f(a) < y < f(b)$
ou $f(a) > y > f(b)$. Então existe $c \in (a,b)$
tal que $f(c) = y$.
\end{thm}

O teorema afirma que uma função contínua em um intervalo fechado
assume todos os valores entre os valores nos extremos.

\begin{proof}
Se $f(a) < y < f(b)$, definimos $g(x) \coloneqq f(x)-y$. Então
$g(a) < 0$ e $g(b) > 0$, e aplicamos o \lemmaref{IVT:lemma}
a $g$ para encontrar $c$. Se $g(c) = 0$, então $f(c) = y$.

De modo semelhante, se $f(a) > y > f(b)$, definimos
$g(x) \coloneqq y-f(x)$. Novamente, $g(a) < 0$ e $g(b) > 0$,
e aplicamos o \lemmaref{IVT:lemma} para encontrar $c$.
Como antes, se $g(c) = 0$, então $f(c) = y$.
\end{proof}

Se uma função é contínua, sua restrição a um subconjunto também é contínua;
se $f \colon S \to \R$ é contínua e $[a,b] \subset S$, então
$f|_{[a,b]}$ também é contínua. Em geral, aplicamos o teorema a uma função
contínua em algum conjunto maior $S$, mas restringimos nossa atenção a um
intervalo.

A demonstração do lema nos diz como encontrar a raiz $c$. Essa prova não é útil
apenas para nós, matemáticos puros: ela também fornece uma ideia útil na
matemática aplicada, chamada de \emph{\myindex{método da bisseção}}.

\begin{example}[Método da bisseção] %\index{bisection method}
O polinômio $f(x) \coloneqq x^3-2x^2+x-1$ tem uma raiz real em $(1,2)$.
Basta observar que $f(1) = -1$ e $f(2) = 1$. Portanto, existe um ponto
$c \in (1,2)$ tal que $f(c) = 0$. Para encontrar uma aproximação melhor
da raiz, seguimos a demonstração do \lemmaref{IVT:lemma}.
Calculamos $f(1.5) = -0.625$. Portanto, há uma raiz do polinômio em $(1.5,2)$.
Em seguida, examinamos 1.75 e observamos que $f(1.75) \approx -0.016$.
Logo, há uma raiz de $f$ em $(1.75,2)$. Depois, calculamos
$f(1.875) \approx 0.44$, portanto há uma raiz em $(1.75,1.875)$.
Repetimos esse procedimento até obter precisão suficiente. De fato,
a raiz é $c \approx 1.7549$.
\end{example}

Essa técnica é o método mais simples para encontrar raízes de polinômios,
um problema comum em matemática aplicada. Em geral, é difícil encontrar raízes
com rapidez, precisão e de modo automático.
Há outros métodos mais rápidos para encontrar raízes de polinômios,
como o método de Newton. Uma vantagem do método acima é sua simplicidade.
Assim que encontramos um intervalo ao qual se aplica o teorema do valor
intermediário, temos a garantia de encontrar uma raiz com a precisão desejada
em um número finito de passos. Além disso, o método da bisseção encontra raízes
de qualquer função contínua, e não apenas de um polinômio.

O teorema garante um $c$ tal que $f(c) = y$, mas pode haver outras raízes
da equação $f(c) = y$. Se seguirmos o procedimento da demonstração,
teremos a garantia de encontrar aproximações para uma dessas raízes.
Para encontrar quaisquer outras raízes, precisamos nos esforçar mais.

\medskip

Polinômios de grau par podem não ter raízes reais.
Não existe número real $x$ tal que $x^2+1 = 0$. Por outro lado,
todo polinômio de grau ímpar tem pelo menos uma raiz real.

\begin{prop}
Seja $f(x)$ um polinômio de grau ímpar. Então $f$ tem uma raiz real.
\end{prop}

\begin{proof}
Suponha que $f$ seja um polinômio de grau ímpar $d$. Escrevemos
\begin{equation*}
f(x) = a_d x^d + a_{d-1} x^{d-1} + \cdots + a_1 x + a_0 ,
\end{equation*}
onde $a_d \neq 0$. Dividimos por $a_d$ para obter um
\emph{\myindex{polinômio mônico}}\footnote{A palavra \emph{mônico} significa que
o coeficiente de $x^d$ é 1.}
\begin{equation*}
g(x) \coloneqq x^d + b_{d-1} x^{d-1} + \cdots + b_1 x + b_0 ,
\end{equation*}
onde $b_k = \nicefrac{a_k}{a_d}$.
Vamos mostrar que $g(n)$ é positivo para algum $n \in \N$ suficientemente
grande. Primeiro, comparamos o termo de ordem mais alta com os demais:
\begin{equation*}
\begin{split}
\abs{\frac{b_{d-1} n^{d-1} + \cdots + b_1 n + b_0}{n^d}}
& =
\frac{\sabs{b_{d-1} n^{d-1} + \cdots + b_1 n + b_0}}{n^d}
\\
& \leq
\frac{\sabs{b_{d-1}} n^{d-1} + \cdots + \sabs{b_1} n + \sabs{b_0}}{n^d}
\\
& \leq
\frac{\sabs{b_{d-1}} n^{d-1} + \cdots + \sabs{b_1} n^{d-1} + \sabs{b_0} n^{d-1}}{n^d}
\\
& =
\frac{n^{d-1}\bigl(\sabs{b_{d-1}} + \cdots + \sabs{b_1} + \sabs{b_0}\bigr)}{n^d}
\\
& =
\frac{1}{n}
\bigl(\sabs{b_{d-1}} + \cdots + \sabs{b_1} + \sabs{b_0}\bigr) .
\end{split}
\end{equation*}
Portanto,
\begin{equation*}
\lim_{n\to\infty} \frac{b_{d-1} n^{d-1} + \cdots + b_1 n + b_0}{n^d}
= 0 .
\end{equation*}
Logo, existe $M \in \N$ tal que
\begin{equation*}
\abs{\frac{b_{d-1} M^{d-1} + \cdots + b_1 M + b_0}{M^d}} < 1 ,
\end{equation*}
o que implica
\begin{equation*}
-(b_{d-1} M^{d-1} + \cdots + b_1 M + b_0) < M^d .
\end{equation*}
Portanto, $g(M) > 0$.

Em seguida, consideramos $g(-n)$, para $n \in \N$. Por um argumento
semelhante, existe $K \in \N$ tal que
$b_{d-1} {(-K)}^{d-1} + \cdots + b_1 (-K) + b_0 < K^d$
e, portanto, $g(-K) < 0$ (veja o \exerciseref{exercise:odddegnegativeK}).
Na demonstração, certifique-se de usar o fato de que $d$ é ímpar.
Em particular, se $d$ é ímpar, então ${(-n)}^d = -(n^d)$.

Aplicamos o teorema do valor intermediário para encontrar
$c \in (-K,M)$ tal que $g(c) = 0$. Como $g(x) = \frac{f(x)}{a_d}$,
temos $f(c) = 0$, e a demonstração está concluída.
\end{proof}

\begin{example}
Você talvez se lembre de quanto nos esforçamos em
\exampleref{example:sqrt2} para mostrar que $\sqrt{2}$ existe.
Com o teorema de Bolzano, podemos provar sem esforço a existência de uma
raiz $k$-ésima de qualquer número positivo $y > 0$.
Afirmamos que, para todo $k \in \N$ e todo $y > 0$,
existe um número $x > 0$ tal que $x^k = y$.

Demonstração: Se $y=1$, isso é claro; portanto, suponhamos $y\neq 1$.
De modo semelhante, suponhamos $k \geq 2$.
Seja $f(x) \coloneqq x^k - y$. Observamos que $f(0) = -y < 0$.
Se $y < 1$, então $f(1) = 1^k -y > 0$. Se $y > 1$,
então $f(y) = y^k-y = y(y^{k-1}-1) > 0$.
Em qualquer dos casos, aplicamos o teorema de Bolzano para encontrar
$x > 0$ tal que $f(x) = 0$, ou, em outras palavras, $x^k = y$.
\end{example}

\begin{example}
Curiosamente, existem funções descontínuas com a propriedade do valor
intermediário. A função
\begin{equation*}
f(x) \coloneqq
\begin{cases}
\sin(\nicefrac{1}{x}) & \text{se } x \neq 0, \\
0 & \text{se } x=0,
\end{cases}
\end{equation*}
não é contínua em $0$; no entanto, $f$ tem a propriedade do valor intermediário:
sempre que $a < b$ e $y$ é tal que $f(a) < y < f(b)$
ou $f(a) > y > f(b)$,
existe $c \in (a,b)$ tal que $f(c) = y$.
Veja a \figureref{figsin1x} para um gráfico de $\sin(\nicefrac{1}{x})$.
A demonstração fica como exercício
(\exerciseref{exercise:meanvaluepropsin1x}).
\end{example}

O teorema do valor intermediário afirma que, se $f \colon [a,b] \to \R$ é
contínua, então $f\bigl([a,b]\bigr)$ contém todos os valores entre $f(a)$ e
$f(b)$. Na verdade, vale algo mais. Combinando todos os resultados desta seção,
podemos provar o seguinte corolário útil, cuja demonstração fica como exercício.
Dica: veja a \figureref{fig:imageinterval} e observe quais são os extremos
do intervalo imagem.

\begin{cor} \label{cor:imageofinterval}
Se $f \colon [a,b] \to \R$ é contínua, então a imagem direta
$f\bigl([a,b]\bigr)$ é um intervalo fechado e limitado ou um único número.
\end{cor}

\begin{myfigureht}
\myincludepdft{figimageinterval}{%
Gráfico de uma função y=f(x) no intervalo de a a b, assinalado por
um traço mais grosso. No eixo vertical, o conjunto de todos os valores
assumidos pela função está marcado como f([a,b]).}
\caption{A imagem de uma função contínua $f \colon [a,b] \to \R$.\label{fig:imageinterval}}
\end{myfigureht}

\subsection{Exercises}

\begin{exercise}
Find an example of a discontinuous function $f \colon [0,1] \to \R$
where the conclusion of the intermediate value theorem fails.
\end{exercise}

\begin{exercise}
Find an example of a \emph{bounded} discontinuous function $f \colon [0,1]
\to \R$ that has neither an absolute minimum nor an absolute maximum.
\end{exercise}

\begin{exercise}
Let $f \colon (0,1) \to \R$ be a continuous function such that
$\displaystyle \lim_{x\to 0} f(x) =
\displaystyle \lim_{x\to 1} f(x) = 0$.  Show that
$f$ achieves either an absolute minimum or an absolute maximum on $(0,1)$
(but perhaps not both).
\end{exercise}

\begin{exercise} \label{exercise:meanvaluepropsin1x}
Let
\begin{equation*}
f(x) \coloneqq
\begin{cases}
\sin(\nicefrac{1}{x}) & \text{if } x \neq 0, \\
0 & \text{if } x=0.
\end{cases}
\end{equation*}
Show that $f$ has the intermediate value property.
That is, whenever $a < b$, if there exists a $y$ such that $f(a) < y < f(b)$
or $f(a) > y > f(b)$, then
there exists a $c \in (a,b)$ such that $f(c) = y$.
\end{exercise}

\begin{exercise} \label{exercise:odddegnegativeK}
Suppose $g(x)$ is a monic polynomial of odd degree $d$, that is,
\begin{equation*}
g(x) = x^d + b_{d-1} x^{d-1} + \cdots + b_1 x + b_0 ,
\end{equation*}
for some real numbers $b_{0}, b_1, \ldots, b_{d-1}$.  Show that there exists
a $K \in \N$ such that $g(-K) < 0$.  Hint: Make sure to use the fact that
$d$ is odd.  You will have to use that ${(-n)}^d = -(n^d)$.
\end{exercise}

\begin{exercise}
Suppose $g(x)$ is a monic polynomial of positive even degree $d$, that is,
\begin{equation*}
g(x) = x^d + b_{d-1} x^{d-1} + \cdots + b_1 x + b_0 ,
\end{equation*}
for some real numbers $b_{0}, b_1, \ldots, b_{d-1}$.  Suppose 
$g(0) < 0$.  Show that $g$ has at least two distinct real roots.
\end{exercise}

\begin{exercise}
Prove \corref{cor:imageofinterval}:
Suppose $f \colon [a,b] \to \R$ is a continuous function.  Prove
that the direct image $f\bigl([a,b]\bigr)$ is a closed and bounded interval or
a single number.
\end{exercise}

\begin{exercise}
Suppose $f \colon \R \to \R$ is continuous and periodic with period
$P > 0$.  That is, $f(x+P) = f(x)$ for all $x \in \R$.  Show that $f$
achieves an absolute minimum and an absolute maximum.
\end{exercise}

\begin{exercise}[Challenging]
Suppose $f(x)$ is a bounded polynomial,
in other words, there is an $M$ such that $\babs{f(x)} \leq M$
for all $x \in \R$.  Prove that $f$ must be a constant.
\end{exercise}

\begin{exercise}
Suppose $f \colon [0,1] \to [0,1]$ is continuous.  Show that $f$
has a fixed point, in other words, show that there exists an $x \in [0,1]$ such that
$f(x) = x$.
\end{exercise}

\begin{exercise}
Find an example of a continuous bounded function $f \colon \R \to \R$ that
achieves neither an absolute minimum nor an absolute maximum on $\R$.
\end{exercise}

\begin{exercise}
Suppose $f \colon \R \to \R$ is continuous such that
$x \leq f(x) \leq x+1$ for all $x \in \R$.  Find $f(\R)$.
\end{exercise}

\begin{exercise}
True/False, prove or find a counterexample.  If $f \colon \R \to
\R$ is a continuous function such that $f|_{\Z}$ is bounded, then $f$
is bounded.
\end{exercise}

\begin{exercise}
Suppose $f \colon [0,1] \to (0,1)$ is a bijection.  Prove that $f$ is not
continuous.
\end{exercise}

\begin{exercise}
Suppose $f \colon \R \to \R$ is continuous.
\begin{enumerate}[a)]
\item
Prove that if there is a $c$ such that $f(c)f(-c) < 0$,
then there is a $d \in \R$ such that $f(d) = 0$.
\item
Find a continuous function $f$ such that
$f(\R) = \R$, but $f(x)f(-x) \geq 0$ for all $x \in \R$.
\end{enumerate}
\end{exercise}

\begin{exercise}
Suppose $g(x)$ is a monic polynomial of even degree $d$, that is,
\begin{equation*}
g(x) = x^d + b_{d-1} x^{d-1} + \cdots + b_1 x + b_0 ,
\end{equation*}
for some real numbers $b_{0}, b_1, \ldots, b_{d-1}$.
Show that $g$ achieves an absolute minimum on $\R$.
\end{exercise}

\begin{exercise}
Suppose $f(x)$ is a polynomial of degree $d$ and 
$f(\R) = \R$.  Show that $d$ is odd.
\end{exercise}

%%%%%%%%%%%%%%%%%%%%%%%%%%%%%%%%%%%%%%%%%%%%%%%%%%%%%%%%%%%%%%%%%%%%%%%%%%%%%%

\sectionnewpage
\section{Uniform continuity}
\label{sec:unifcont}

%mbxINTROSUBSECTION

\sectionnotes{1,5--2 aulas (a extensão contínua pode ser opcional)}

\subsection{Uniform continuity}

Enfatizamos bastante que o $\delta$ na definição de continuidade dependia
do ponto $c$. Há situações em que é vantajoso poder escolher um $\delta$
independente de qualquer ponto; por isso, damos um nome a esse conceito.

\begin{defn}
Seja $S \subset \R$ e seja $f \colon S \to \R$ uma função.
Suponha que, para todo $\epsilon > 0$, exista $\delta > 0$
tal que, sempre que $x,c \in S$ e
$\sabs{x-c} < \delta$, tenhamos $\babs{f(x)-f(c)} < \epsilon$.
Dizemos então que $f$ é \emph{\myindex{uniformemente contínua}}.
\end{defn}

Uma função uniformemente contínua precisa ser contínua.
A única diferença entre as definições é que, na continuidade uniforme,
dado $\epsilon > 0$, escolhemos um $\delta > 0$ que funciona para todo
$c \in S$. Isto é, $\delta$ não pode mais depender de $c$;
depende apenas de $\epsilon$. Agora, o domínio da função faz diferença.
Uma função que não é uniformemente contínua em um conjunto maior pode ser
uniformemente contínua quando restrita a um conjunto menor.
Observe que, nesta definição, $x$ e $c$ não recebem tratamentos diferentes.

\begin{example}
A função $f \colon [0,1] \to \R$ definida por $f(x) \coloneqq x^2$
é uniformemente contínua.

Demonstração: Observe que $0 \leq x,c \leq 1$. Então
\begin{equation*}
\sabs{x^2-c^2} = \sabs{x+c}\sabs{x-c}
\leq \bigl(\sabs{x}+\sabs{c}\bigr)\sabs{x-c}
\leq (1+1)\sabs{x-c} .
\end{equation*}
Portanto, dado $\epsilon > 0$, seja $\delta \coloneqq \nicefrac{\epsilon}{2}$.
Se $\sabs{x-c} < \delta$, então
$\sabs{x^2-c^2} \leq 2\sabs{x-c} < \epsilon$.

\medskip

Por outro lado, a função $g \colon \R \to \R$ definida por $g(x) \coloneqq x^2$
não é uniformemente contínua.

Demonstração: Suponha que seja uniformemente contínua.
Então, para todo $\epsilon > 0$, existiria $\delta > 0$ tal que,
se $\sabs{x-c} < \delta$, então $\sabs{x^2-c^2} < \epsilon$.
Tome $x > 0$ e seja
$c \coloneqq x+\nicefrac{\delta}{2}$. Escrevemos
\begin{equation*}
\epsilon >
\sabs{x^2-c^2} = \sabs{x+c}\sabs{x-c}
=
(2x+\nicefrac{\delta}{2})\nicefrac{\delta}{2}
\geq
\delta x .
\end{equation*}
Portanto, $x < \nicefrac{\epsilon}{\delta}$ para todo $x > 0$,
o que é uma contradição.
\end{example}


\begin{example}
A função $f \colon (0,1) \to \R$ definida por
$f(x) \coloneqq \nicefrac{1}{x}$ não é uniformemente contínua.

Demonstração: Dado $\epsilon > 0$, a desigualdade
$\epsilon > \abs{\nicefrac{1}{x}-\nicefrac{1}{y}}$ vale se, e somente se,
\begin{equation*}
\epsilon >
\abs{\nicefrac{1}{x}-\nicefrac{1}{y}}
=
\frac{\sabs{y-x}}{\sabs{xy}}
=
\frac{\sabs{y-x}}{xy} ,
\end{equation*}
ou
\begin{equation*}
\sabs{x-y} < xy\epsilon .
\end{equation*}
Suponha que $\epsilon < 1$.
Queremos verificar se um $\delta > 0$ pequeno funcionaria.
Se $x \in (0,1)$ e
$y = x+\nicefrac{\delta}{2} \in (0,1)$,
então $\sabs{x-y} = \nicefrac{\delta}{2} < \delta$.
Substituímos $y$ na desigualdade acima e obtemos
$\nicefrac{\delta}{2} < x \bigl(x+\nicefrac{\delta}{2}\bigr)\epsilon < x$.
Se a definição de continuidade uniforme fosse satisfeita,
a desigualdade $\nicefrac{\delta}{2} < x$ valeria para todo $x > 0$
suficientemente pequeno. Mas isso implica que $\delta \leq 0$.
Portanto, nenhum $\delta > 0$ único funciona para todos os pontos.
\end{example}

Os exemplos mostram que, se $f$ está definida em um intervalo que não é
fechado ou não é limitado, $f$ pode ser contínua, mas não uniformemente contínua.
Para um intervalo fechado e limitado $[a,b]$, porém, podemos fazer a seguinte
afirmação.

\begin{thm} \label{unifcont:thm}
Seja $f \colon [a,b] \to \R$ uma função contínua. Então $f$
é uniformemente contínua.
\end{thm}

\begin{proof}
Provaremos a afirmação por contraposição.
Suponha que $f$ não seja uniformemente contínua. Provaremos
que existe algum $c \in [a,b]$ em que $f$ não é contínua.
Neguemos a definição de continuidade uniforme:
existe $\epsilon > 0$
tal que, para todo $\delta > 0$, existem pontos $x,y$ em $[a,b]$ com
$\sabs{x-y} < \delta$ e $\babs{f(x)-f(y)} \geq \epsilon$.

Assim, para o $\epsilon > 0$ acima,
encontramos sequências $\{ x_n \}_{n=1}^\infty$ e $\{ y_n \}_{n=1}^\infty$
tais que $\sabs{x_n-y_n} < \nicefrac{1}{n}$ e
$\babs{f(x_n)-f(y_n)} \geq \epsilon$. Pelo
\hyperref[thm:bwseq]{teorema de Bolzano--Weierstrass},
existe uma subsequência convergente
$\{ x_{n_k} \}_{k=1}^\infty$. Seja
$c \coloneqq \lim_{k\to\infty} x_{n_k}$.
Como $a \leq x_{n_k} \leq b$ para todo $k$, temos $a \leq c \leq b$. Estimamos:
\begin{equation*}
\sabs{y_{n_k} - c} =
\sabs{y_{n_k} - x_{n_k} + x_{n_k} - c} \leq
\sabs{y_{n_k} - x_{n_k}}
+
\sabs{x_{n_k}-c}
<
\nicefrac{1}{n_k}
+
\sabs{x_{n_k}-c} .
\end{equation*}
Como $\nicefrac{1}{n_k}$ e $\sabs{x_{n_k}-c}$ tendem ambos a zero quando
$k$ tende ao infinito, a sequência $\{ y_{n_k} \}_{k=1}^\infty$ converge,
e seu limite é $c$. Agora mostraremos que $f$ não é contínua em $c$.
Estimamos:
\begin{equation*}
\begin{split}
\babs{f(x_{n_k}) - f(c)} & =
\babs{f(x_{n_k}) - f(y_{n_k}) + f(y_{n_k}) - f(c)} \\
& \geq
\babs{f(x_{n_k}) - f(y_{n_k})} - \babs{f(y_{n_k}) - f(c)} \\
& \geq
\epsilon - \babs{f(y_{n_k})-f(c)} .
\end{split}
\end{equation*}
Ou, em outras palavras,
\begin{equation*}
\babs{f(x_{n_k})-f(c)}
+
\babs{f(y_{n_k})-f(c)} \geq
\epsilon .
\end{equation*}
Pelo menos uma das sequências $\bigl\{ f(x_{n_k}) \bigr\}_{k=1}^\infty$ ou
$\bigl\{ f(y_{n_k}) \bigr\}_{k=1}^\infty$ não pode convergir para $f(c)$.
Caso contrário, o lado esquerdo da desigualdade tenderia a zero, enquanto
o lado direito é positivo. Portanto, $f$ não pode ser contínua em $c$.
\end{proof}

Como antes, observe o ponto essencial da demonstração: podemos aplicar
Bolzano--Weierstrass porque o intervalo $[a,b]$ é limitado, e o limite da
subsequência pertence novamente a $[a,b]$ porque o intervalo é fechado.

\subsection{Continuous extension}

Funções uniformemente contínuas em intervalos abertos podem ser estendidas
continuamente aos extremos.
A ideia central é o lema a seguir, que também tem muitos outros usos.
Ele afirma que funções uniformemente contínuas preservam sequências de Cauchy.
A nova questão aqui é que, para uma sequência de Cauchy,
o limite pode não pertencer ao domínio da função.

\begin{lemma} \label{unifcauchycauchy:lemma}
Seja $S \subset \R$ e seja $f \colon S \to \R$ uma função uniformemente contínua.
Seja $\{ x_n \}_{n=1}^\infty$ uma sequência de Cauchy em $S$.
Então $\bigl\{ f(x_n) \bigr\}_{n=1}^\infty$ é de Cauchy.
\end{lemma}

\begin{proof}
Seja $\epsilon > 0$ dado. Existe $\delta > 0$ tal que
$\babs{f(x)-f(y)} < \epsilon$ sempre que $x,y \in S$ e
$\sabs{x-y} < \delta$. Encontre $M \in \N$ tal que, para todos $n,k \geq M$,
tenhamos $\sabs{x_n-x_k} < \delta$.
Então, para todos $n,k \geq M$, temos $\babs{f(x_n)-f(x_k)} < \epsilon$.
\end{proof}

Uma aplicação do lema acima é o seguinte resultado de extensão. Ele afirma que
uma função em um intervalo aberto é uniformemente contínua se, e somente se,
pode ser estendida a uma função contínua no intervalo fechado.

\begin{prop} \label{context:prop}
Uma função $f \colon (a,b) \to \R$ é uniformemente contínua se, e somente se,
os limites
\begin{equation*}
L_a \coloneqq \lim_{x \to a} f(x) \qquad \text{e} \qquad
L_b \coloneqq \lim_{x \to b} f(x)
\end{equation*}
existem e a função $\widetilde{f} \colon [a,b] \to \R$,
definida por
\begin{equation*}
\widetilde{f}(x) \coloneqq
\begin{cases}
f(x) & \text{se } x \in (a,b), \\
L_a & \text{se } x = a, \\
L_b & \text{se } x = b
\end{cases}
\end{equation*}
é contínua.
\end{prop}

\begin{proof}
Uma das implicações é imediata. Se $\widetilde{f}$ é contínua,
então é uniformemente contínua pelo \thmref{unifcont:thm}. Como $f$ é a
restrição de $\widetilde{f}$ a $(a,b)$, $f$ também é uniformemente contínua
(exercício).

Suponha agora que $f$ seja uniformemente contínua. Primeiro, precisamos mostrar
que os limites $L_a$ e $L_b$ existem. Vamos nos concentrar em $L_a$.
Tome $\{ x_n \}_{n=1}^\infty$ em $(a,b)$ tal que
$\lim_{n\to\infty} x_n = a$.
A sequência $\{ x_n \}_{n=1}^\infty$ é de Cauchy; portanto, pelo
\lemmaref{unifcauchycauchy:lemma},
a sequência $\bigl\{ f(x_n) \bigr\}_{n=1}^\infty$ é de Cauchy e, assim, converge.
Seja $L_1 \coloneqq \lim_{n\to\infty} f(x_n)$. Tome outra sequência
$\{ y_n \}_{n=1}^\infty$ em $(a,b)$ tal que $\lim_{n\to\infty} y_n = a$.
Pelo mesmo raciocínio, obtemos
$L_2 \coloneqq \lim_{n\to\infty} f(y_n)$. Se mostrarmos que $L_1 = L_2$,
então o limite $L_a = \lim_{x\to a} f(x)$ existe. Seja $\epsilon > 0$ dado.
Encontre $\delta > 0$ tal que $\sabs{x-y} < \delta$ implique
$\babs{f(x)-f(y)} < \nicefrac{\epsilon}{3}$. Encontre $M \in \N$ tal que, para
$n \geq M$, tenhamos $\sabs{a-x_n} < \nicefrac{\delta}{2}$,
$\sabs{a-y_n} < \nicefrac{\delta}{2}$,
$\babs{f(x_n)-L_1} < \nicefrac{\epsilon}{3}$ e
$\babs{f(y_n)-L_2} < \nicefrac{\epsilon}{3}$. Então, para $n \geq M$,
\begin{equation*}
\sabs{x_n-y_n} =
\sabs{x_n-a+a-y_n} \leq
\sabs{x_n-a}+\sabs{a-y_n} < \nicefrac{\delta}{2} + \nicefrac{\delta}{2} =
\delta.
\end{equation*}
Logo,
\begin{equation*}
\begin{split}
\sabs{L_1-L_2} &=
\babs{L_1-f(x_n)+f(x_n)-f(y_n)+f(y_n)-L_2} \\
& \leq
\babs{L_1-f(x_n)}+\babs{f(x_n)-f(y_n)}+\babs{f(y_n)-L_2} \\
& <
\nicefrac{\epsilon}{3} + \nicefrac{\epsilon}{3} + \nicefrac{\epsilon}{3}
=
\epsilon .
\end{split}
\end{equation*}
Portanto, $L_1 = L_2$.
Assim, $L_a$ existe. Mostrar que $L_b$ existe fica como exercício.

Se $L_a = \lim_{x\to a} f(x)$ existe, então $\lim_{x\to a} \widetilde{f}(x)$
existe e é igual a $L_a$ (veja a \propref{prop:limrest}).
De modo semelhante, o mesmo vale para $L_b$.
Logo, $\widetilde{f}$ é contínua em $a$ e em $b$.
E, como $f$ é contínua em $c \in (a,b)$,
concluímos que $\widetilde{f}$ é contínua em $c \in (a,b)$
(\propref{prop:limrest}, novamente).
\end{proof}

Uma aplicação típica desta proposição (junto com
\propref{prop:onesidedlimits}) é a seguinte.
Se $f \colon (-1,0) \cup (0,1) \to \R$ é uniformemente contínua,
então $\lim_{x\to 0} f(x)$ existe e a função tem uma
\emph{\myindex{singularidade removível}}.
Isto é, podemos estender a função a uma função contínua em $(-1,1)$.


\subsection{Lipschitz continuous functions}

\begin{defn}
Uma função $f \colon S \to \R$
é \emph{\myindex{Lipschitz contínua}}\index{função!Lipschitz}%
\footnote{Nomeada em homenagem ao matemático alemão
\href{https://en.wikipedia.org/wiki/Rudolf_Lipschitz}{Rudolf Otto Sigismund Lipschitz}
(1832--1903).} se existe $K \in \R$ tal que
\begin{equation*}
\babs{f(x)-f(y)} \leq K \sabs{x-y}
\qquad \text{para todo } x \text{ and } y \text{ in } S.
\end{equation*}
\end{defn}

Uma grande classe de funções é Lipschitz contínua. Cuidado: assim como no
caso das funções uniformemente contínuas, o domínio da função é importante.
Veja os exemplos abaixo e os exercícios. Primeiro, justificamos o uso da
palavra \emph{contínua}.

\begin{prop}
Uma função Lipschitz contínua é uniformemente contínua.
\end{prop}

\begin{proof}
Seja $f \colon S \to \R$ uma função e seja $K$ uma constante tal que
$\babs{f(x)-f(y)} \leq K \sabs{x-y}$
para todos $x, y$ em $S$.
Seja $\epsilon > 0$ dado. Tome $\delta \coloneqq
\nicefrac{\epsilon}{K}$.
Para todos $x$ e $y$ em $S$ tais que
$\sabs{x-y} < \delta$,
\begin{equation*}
\babs{f(x)-f(y)} \leq K \sabs{x-y} < K \delta = K \frac{\epsilon}{K} =
\epsilon .
\end{equation*}
Portanto, $f$ é uniformemente contínua.
\end{proof}

Interpretamos geometricamente a continuidade de Lipschitz. Seja $f$ uma função
Lipschitz contínua com alguma constante $K$. Reescrevemos a desigualdade
para dizer que, para $x \neq y$, temos
\begin{equation*}
\abs{\frac{f(x)-f(y)}{x-y}} \leq K .
\end{equation*}
A quantidade $\frac{f(x)-f(y)}{x-y}$ é a inclinação da reta que passa pelos
pontos $\bigl(x,f(x)\bigr)$ e $\bigl(y,f(y)\bigr)$, isto é, uma
\emph{\myindex{reta secante}}. Portanto, $f$ é Lipschitz contínua se, e
somente se, toda reta que intersecta o gráfico de $f$ em pelo menos dois
pontos distintos tem inclinação, em valor absoluto, menor ou igual a $K$.
Veja a \figureref{fig:lipschitz}.
\begin{myfigureht}
\myincludegraphics{lipschitzfig}{%
Gráfico de uma função com alguns vértices. Dois pontos no eixo horizontal
estão marcados com x e y. Uma reta que passa pelos pontos correspondentes
no gráfico está desenhada, e sua inclinação está indicada como
f de x menos f de y, tudo dividido por x menos y.}
\caption{A inclinação de uma reta secante.
Uma função é Lipschitz se %$\abs{\text{slope}} =
$\abs{\frac{f(x)-f(y)}{x-y}} \leq K$ para todos $x$ e $y$. \label{fig:lipschitz}}
\end{myfigureht}

\begin{example}
As funções $\sin(x)$ e $\cos(x)$ são Lipschitz contínuas.
No \exampleref{sincos:example}, vimos as seguintes duas desigualdades:
\begin{equation*}
\babs{\sin(x)-\sin(y)}
\leq \sabs{x-y}
\qquad \text{e} \qquad
\babs{\cos(x)-\cos(y)}
\leq \sabs{x-y} .
\end{equation*}

Logo, seno e cosseno são Lipschitz contínuas com $K=1$.
\end{example}

\begin{example}
A função $f \colon [1,\infty) \to \R$ definida por $f(x) \coloneqq \sqrt{x}$
é Lipschitz contínua. Demonstração:
\begin{equation*}
\abs{\sqrt{x}-\sqrt{y}} =
\abs{\frac{x-y}{\sqrt{x}+\sqrt{y}}}
=
\frac{\sabs{x-y}}{\sqrt{x}+\sqrt{y}} .
\end{equation*}
Como $x \geq 1$ e $y \geq 1$, vemos que
$\frac{1}{\sqrt{x}+\sqrt{y}} \leq \frac{1}{2}$. Portanto,
\begin{equation*}
\abs{\sqrt{x}-\sqrt{y}} =
\abs{\frac{x-y}{\sqrt{x}+\sqrt{y}}}
\leq
\frac{1}{2}
\sabs{x-y}.
\end{equation*}

Por outro lado, a função $g \colon [0,\infty) \to \R$ definida por
$g(x) \coloneqq \sqrt{x}$ não é Lipschitz contínua. Demonstração:
Suponha que exista $K$ tal que, para todos $x,y \in [0,\infty)$,
\begin{equation*}
\abs{\sqrt{x}-\sqrt{y}}
\leq
K \sabs{x-y} ,
\end{equation*}
Concluímos que $K > 0$, pois $K < 0$ é impossível,
e $K=0$ significaria que $f$ é constante.
Tomamos $y=0$ para obter
$\sqrt{x} \leq Kx$.
Para $x > 0$, obtemos
$\nicefrac{1}{K} \leq \sqrt{x}$ ou $\nicefrac{1}{K^2} \leq x$.
Isso não pode ser verdade para todo $x > 0$. Assim, não existe $K$,
e $g$ não é Lipschitz contínua. Veja a \figureref{fig:sqrtgraph} e observe
como as retas secantes ficariam cada vez mais verticais ao nos aproximarmos
de $x=0$.
\begin{myfigureht}
\myincludegraphics{sqrtgraph}{%
Gráfico de uma função que começa subindo verticalmente no extremo esquerdo
e logo se curva, tornando-se muito mais horizontal no extremo direito.
Estão desenhadas retas secantes que passam pelo extremo esquerdo e por uma
sequência de pontos no gráfico que se aproximam cada vez mais desse extremo.
Observamos que as retas secantes ficam cada vez mais verticais à medida que
o segundo ponto da reta secante se aproxima do extremo esquerdo.}
\caption{Gráfico de $\sqrt{x}$ e algumas retas secantes que passam por
$(0,0)$ e $(x,\sqrt{x})$.\label{fig:sqrtgraph}}
\end{myfigureht}

O último exemplo mostra uma função $g$ uniformemente contínua,
mas não Lipschitz contínua.
Para ver que $\sqrt{x}$ é uniformemente contínua como função em $[0,\infty)$,
observe que ela é uniformemente contínua quando restrita a $[0,1]$ pelo
\thmref{unifcont:thm}. Ela também é Lipschitz (e, portanto, uniformemente
contínua) quando restrita a $[1,\infty)$.
Não é difícil mostrar (exercício) que isso significa que $\sqrt{x}$ é uma
função uniformemente contínua em $[0,\infty)$.
\end{example}

\subsection{Exercises}

\begin{exercise}
Seja $f \colon S \to \R$ uniformemente contínua. Seja $A \subset S$.
Então a restrição $f|_A$ é uniformemente contínua.
\end{exercise}

\begin{exercise}
Seja $f \colon (a,b) \to \R$ uma função uniformemente contínua.
Complete a demonstração da \propref{context:prop}, mostrando que
o limite $\lim\limits_{x \to b} f(x)$ existe.
\end{exercise}

\begin{exercise}
Mostre que $f \colon (c,\infty) \to \R$, para algum $c > 0$,
definida por $f(x) \coloneqq \nicefrac{1}{x}$ é Lipschitz contínua.
\end{exercise}

\begin{exercise}
Mostre que $f \colon (0,\infty) \to \R$
definida por $f(x) \coloneqq \nicefrac{1}{x}$ não é Lipschitz contínua.
\end{exercise}

\begin{samepage}
\begin{exercise}
Sejam $A,B$ intervalos.
Sejam $f \colon A \to \R$ e $g \colon B \to \R$ funções uniformemente contínuas
tais que $f(x) = g(x)$ para $x \in A \cap B$. Defina
a função $h \colon A \cup B \to \R$ por $h(x) \coloneqq f(x)$ se
$x \in A$ e $h(x) \coloneqq g(x)$ se $x \in B \setminus A$.
\begin{enumerate}[a)]
\item
Prove que, se $A \cap B \neq \emptyset$, então $h$ é uniformemente contínua.
\item
Encontre um exemplo em que $A \cap B = \emptyset$ e $h$ nem sequer
seja contínua.
\end{enumerate}
\end{exercise}
\end{samepage}

\begin{exercise}[Desafiador]
Seja $f \colon \R \to \R$ um polinômio de grau
$d \geq 2$. Mostre que $f$ não é Lipschitz contínua.
\end{exercise}

\begin{exercise}
Seja $f \colon (0,1) \to \R$ uma função contínua e limitada. Mostre que
a função $g(x) \coloneqq x(1-x)f(x)$ é uniformemente contínua.
\end{exercise}

\begin{exercise}
Mostre que $f \colon (0,\infty) \to \R$ definida por
$f(x) \coloneqq \sin(\nicefrac{1}{x})$ não é uniformemente contínua.
\end{exercise}

\begin{exercise}[Desafiador]
Seja $f \colon \Q \to \R$ uma função uniformemente contínua. Mostre que
existe uma função uniformemente contínua $\widetilde{f} \colon \R \to \R$
tal que $f(x) = \widetilde{f}(x)$ para todo $x \in \Q$.
\end{exercise}

\begin{samepage}
\begin{exercise}
\leavevmode
\begin{enumerate}[a)]
\item
Encontre uma função contínua $f \colon (0,1) \to \R$ e uma sequência
$\{ x_n \}_{n=1}^\infty$ em $(0,1)$ que seja de Cauchy, mas tal que
$\bigl\{ f(x_n) \bigr\}_{n=1}^\infty$ não seja de Cauchy.
\item
Prove que, se $f \colon \R \to \R$ é contínua e
$\{ x_n \}_{n=1}^\infty$ é de Cauchy, então
$\bigl\{ f(x_n) \bigr\}_{n=1}^\infty$ é de Cauchy.
\end{enumerate}
\end{exercise}
\end{samepage}

\begin{samepage}
\begin{exercise}
Demonstre:
\begin{enumerate}[a)]
\item
Se $f \colon S \to \R$ e $g \colon S \to \R$ são uniformemente contínuas,
então $h \colon S \to \R$ definida por $h(x) \coloneqq f(x) + g(x)$
é uniformemente contínua.
\item
Se $f \colon S \to \R$ é uniformemente contínua e $a \in \R$,
então $h \colon S \to \R$ definida por $h(x) \coloneqq a f(x)$
é uniformemente contínua.
\end{enumerate}
\end{exercise}
\end{samepage}

\begin{exercise}
Demonstre:
\begin{enumerate}[a)]
\item
Se $f \colon S \to \R$ e $g \colon S \to \R$ são Lipschitz,
então $h \colon S \to \R$ definida por $h(x) \coloneqq f(x) + g(x)$
é Lipschitz.
\item
Se $f \colon S \to \R$ é Lipschitz e $a \in \R$,
então $h \colon S \to \R$ definida por $h(x) \coloneqq a f(x)$
é Lipschitz.
\end{enumerate}
\end{exercise}

\begin{exercise}
\pagebreak[2]
\leavevmode
\begin{enumerate}[a)]
\item
Se $f \colon [0,1] \to \R$ é dada por $f(x) \coloneqq x^m$ para um inteiro
$m \geq 0$, mostre que $f$ é Lipschitz e encontre a melhor constante de
Lipschitz $K$ (a menor, que depende de $m$, é claro).
Dica:
$(x-y)(x^{m-1} + x^{m-2}y + x^{m-3}y^2 + \cdots + x y^{m-2} + y^{m-1}) = x^m
- y^m$ para $m$ positivo.
\item
Usando o exercício anterior, mostre que, se $f \colon [0,1] \to \R$
é um polinômio, isto é,
$f(x) \coloneqq a_m x^m + a_{m-1} x^{m-1} + \cdots + a_0$,
então $f$ é Lipschitz.
\end{enumerate}
\end{exercise}

\begin{exercise}
\pagebreak[2]
Suponha que, para $f \colon [0,1] \to \R$, tenhamos
$\babs{f(x)-f(y)} \leq K \sabs{x-y}$ para todos $x,y$ em $[0,1]$,
e que $f(0) = f(1) = 0$.
Prove que $\babs{f(x)} \leq \nicefrac{K}{2}$ para todo $x \in [0,1]$.
Mostre também, por meio de um exemplo, que $\nicefrac{K}{2}$ é o melhor
limite possível. Isto é, encontre uma função contínua desse tipo
para a qual $\babs{f(x)} = \nicefrac{K}{2}$ para algum $x \in [0,1]$.
\end{exercise}

\begin{exercise}
Suponha que $f \colon \R \to \R$ seja contínua e periódica, com período
$P > 0$. Isto é, $f(x+P) = f(x)$ para todo $x \in \R$. Mostre que $f$
é uniformemente contínua.
\end{exercise}

\begin{exercise}
Suponha que $f \colon S \to \R$ e $g \colon [0,\infty) \to [0,\infty)$
sejam funções, que $g$ seja contínua em $0$, que $g(0) = 0$ e que,
sempre que $x$ e $y$ estiverem em $S$, tenhamos
$\babs{f(x)-f(y)} \leq g\bigl(\sabs{x-y}\bigr)$.
Prove que $f$ é uniformemente contínua.
\end{exercise}

\begin{exercise}
Suponha que $f \colon [a,b] \to \R$ seja uma função tal que, para todo
$c \in [a,b]$, existam $K_c > 0$ e $\epsilon_c > 0$ para os quais
$\babs{f(x)-f(y)} \leq K_c \sabs{x-y}$ para todos $x$ e $y$ em
$(c-\epsilon_c,c+\epsilon_c) \cap [a,b]$. Em outras palavras, $f$ é
\myquote{localmente Lipschitz}.
\begin{enumerate}[a)]
\item
Prove que existe uma constante global $K > 0$ tal que
$\babs{f(x)-f(y)} \leq K \sabs{x-y}$ para todos $x,y$ em $[a,b]$.
\item
Encontre um contraexemplo ao item anterior quando o intervalo é aberto;
isto é, encontre uma função $f \colon (a,b) \to \R$ que seja localmente
Lipschitz, mas não Lipschitz.
\end{enumerate}
\end{exercise}

%%%%%%%%%%%%%%%%%%%%%%%%%%%%%%%%%%%%%%%%%%%%%%%%%%%%%%%%%%%%%%%%%%%%%%%%%%%%%%

\sectionnewpage
\section{Limits at infinity}
\label{sec:limitatinf}

%mbxINTROSUBSECTION

\sectionnotes{less than 1 lecture (optional, can safely be omitted unless
\sectionref{sec:monotonefunc} or
\sectionref{sec:impropriemann} is also covered)}

\subsection{Limits at infinity}

Just as with sequences, a continuous variable can also approach infinity.

\begin{defn}
We say $\infty$ is a \emph{\myindex{cluster point}} of $S \subset \R$ if for every
$M \in \R$, there exists an $x \in S$ such that $x \geq M$.  Similarly,
$- \infty$ is a \emph{cluster point} of $S \subset \R$ if for every
$M \in \R$, there exists an $x \in S$ such that $x \leq M$.

\index{limit!of a function at infinity}%
Let $f \colon S \to \R$ be a function, where 
$\infty$ is a cluster point of $S$.
If there exists an $L \in \R$
with the property
that for every $\epsilon > 0$, there is an $M \in \R$ such that
\begin{equation*}
\babs{f(x) - L} < \epsilon 
\end{equation*}
whenever $x \in S$ and $x \geq M$, then we say $f(x)$
\emph{converges}\index{converges!function} to $L$
as $x$ goes to $\infty$.  We call $L$ a \emph{limit}
and, if unique, write
\glsadd{not:limfuncinf}
\begin{equation*}
\lim_{x \to \infty} f(x) \coloneqq L .
\end{equation*}
Alternatively, we write $f(x) \to L$ as $x \to \infty$.

Similarly, if $-\infty$ is a cluster point of $S$
and
there exists an $L \in \R$
such that for every $\epsilon > 0$, there is an $M \in \R$ such that
\begin{equation*}
\babs{f(x) - L} < \epsilon 
\end{equation*}
whenever $x \in S$ and $x \leq M$, then we say $f(x)$ \emph{converges} to $L$
as $x$ goes to $-\infty$.
Alternatively, we write $f(x) \to L$ as $x \to -\infty$.
We call $L$ a \emph{limit} and, if unique, write
\begin{equation*}
\lim_{x \to -\infty} f(x) \coloneqq L .
\end{equation*}
\end{defn}

The first thing to do, as usual, is to prove that the limit, if it exists,
is unique.
We leave it as an exercise for the reader.

\begin{prop} \label{liminfty:unique}
The limit at $\infty$ or $-\infty$ as defined above is unique if it exists.
\end{prop}

\begin{example}
Let $f(x) \coloneqq \frac{1}{\sabs{x}+1}$.  Then
\begin{equation*}
\lim_{x\to \infty} f(x) = 0 \qquad \text{and} \qquad
\lim_{x\to -\infty} f(x) = 0 .
\end{equation*}

Proof:
Let $\epsilon > 0$ be given.  Find $M > 0$ large enough
so that $\frac{1}{M+1} < \epsilon$.  If
$x \geq M$, then $0 < \frac{1}{\sabs{x}+1} = \frac{1}{x+1} \leq \frac{1}{M+1} < \epsilon$.
The first limit follows.
The proof for $-\infty$ is left to the reader.
\end{example}

\begin{example}
Let $f(x) \coloneqq \sin(\pi x)$.  Then $\lim_{x\to\infty} f(x)$ does not exist.
To prove this fact, note that if $x = 2n+\nicefrac{1}{2}$ for some $n \in
\N$, then $f(x)=1$,
while if $x = 2n+\nicefrac{3}{2}$, then $f(x)=-1$.
Both of these values, $1$ and $-1$, cannot be
within a small $\epsilon$ of a single real number.

Be careful not to confuse continuous limits with limits of sequences.
We could say
\begin{equation*}
\lim_{n \to \infty} \sin(\pi n) = 0, \qquad \text{but} \qquad
\lim_{x \to \infty} \sin(\pi x) \enspace \text{does not exist}.
\end{equation*}
Of course, the notation is ambiguous:  Are we thinking of the
sequence $\bigl\{ \sin (\pi n) \bigr\}_{n=1}^\infty$ or the function $\sin(\pi x)$
of a real variable?  We are simply using the convention
that $n \in \N$, while $x \in \R$.  When the notation is not clear,
it is good to explicitly mention where the variable lives, or what kind
of limit you are using.  If there is a possibility of confusion, one can
write, for example,
\begin{equation*}
\lim_{\substack{n \to \infty\\n \in \N}} \sin(\pi n) .
\end{equation*}
\end{example}

There is a connection between continuous limits and limits of sequences, but we must take all
sequences going to infinity, just as before in \lemmaref{seqflimit:lemma}.

\begin{lemma} \label{seqflimitinf:lemma}
Suppose $f \colon S \to \R$ is a function, $\infty$ is a cluster
point of $S \subset \R$, and $L \in \R$.  Then
\begin{equation*}
\lim_{x\to\infty} f(x) = L
\qquad \text{if and only if} \qquad
\lim_{n\to\infty} f(x_n) = L
\end{equation*}
for all sequences $\{ x_n \}_{n=1}^\infty$ in $S$ such that $\lim\limits_{n\to\infty} x_n = \infty$.
\end{lemma}

The lemma also holds for the limit as $x \to -\infty$.
Its proof is almost identical and
is left as an exercise.

\begin{proof}
First suppose $f(x) \to L$ as $x \to \infty$.
Given an $\epsilon > 0$, there exists an $M$ such that for all $x \in S$
where $x \geq M$,
we have $\babs{f(x)-L} < \epsilon$.
Let $\{ x_n \}_{n=1}^\infty$
be a sequence in $S$ such that $\lim_{n\to\infty} x_n = \infty$.  Then there exists an
$N$ such that
$x_n \geq M$
for all $n \geq N$.
Thus
$\babs{f(x_n)-L} < \epsilon$.

We prove the converse by contrapositive.  Suppose $f(x)$ does
not go to $L$ as $x \to \infty$.
This means that there exists an $\epsilon > 0$
such that for every $n \in \N$, there exists an $x \in S$, $x \geq n$, let
us call it $x_n$, such that $\babs{f(x_n)-L} \geq \epsilon$.
Consider the sequence $\{ x_n \}_{n=1}^\infty$.  Clearly 
$\bigl\{ f(x_n) \bigr\}_{n=1}^\infty$ does not converge to $L$.  It remains to note
that $\lim_{n\to\infty} x_n = \infty$ because $x_n \geq n$ for all $n$.
\end{proof}

Using the lemma, we again translate results about sequential
limits into results about continuous limits as $x$ goes to infinity.  That
is, we have almost immediate analogues of the corollaries
in \sectionref{subseq:sequentiallimits}.  We simply allow 
the cluster point $c$ to be either $\infty$ or $-\infty$, in addition
to a real number.  We leave it to
the student to verify these statements.

\subsection{Infinite limit}

Just as for sequences, it is often convenient to distinguish certain
divergent sequences, and talk about limits being infinite
almost as if the limits existed.

\begin{defn}
\index{infinite limit!of a function}\index{limit!infinite}%
Let $f \colon S \to \R$ be a function and suppose 
$S$ has $\infty$ as a cluster point.
We say $f(x)$
\emph{\myindex{diverges to infinity}} 
as $x$ goes to $\infty$
if for every $N \in \R$
there exists an $M \in \R$ such that
\begin{equation*}
f(x) > N
\end{equation*}
whenever $x \in S$ and $x \geq M$.
We write
\begin{equation*}
\lim_{x \to \infty} f(x) \coloneqq \infty ,
\end{equation*}
or we say that $f(x) \to \infty$ as $x \to \infty$.
\end{defn}

A similar definition can be made for limits as $x \to -\infty$
or as $x \to c$ for a finite $c$.  Also, similar definitions can be
made for limits being $-\infty$.  Stating these definitions is left
as an exercise.
Note that
sometimes the slightly incorrect (in this context)
wording
\emph{\myindex{converges to infinity}} is used.
We can again use sequential limits, and an analogue of 
\lemmaref{seqflimit:lemma} is left as an exercise.

\begin{example}
Let us show that $\lim\limits_{x \to \infty} \frac{1+x^2}{1+x} = \infty$.

Proof: For $x \geq 1$, we have
\begin{equation*}
\frac{1+x^2}{1+x} \geq 
\frac{x^2}{x+x}  = 
\frac{x}{2} .
\end{equation*}
Given $N \in \R$, take $M = \max \{ 2N+1 , 1 \}$.
If $x \geq M$, then $x \geq 1$ and $\nicefrac{x}{2} > N$.
So
\begin{equation*}
\frac{1+x^2}{1+x} \geq 
\frac{x}{2} > N .
\end{equation*}
\end{example}

\subsection{Compositions}

Finally, just as for limits at finite numbers,
limits play nice with compositions given a
continuity assumption.

\begin{prop} \label{prop:inflimcompositions}
Suppose $f \colon A \to B$, $g \colon B \to \R$, $A, B \subset \R$, 
$a \in \R \cup \{ -\infty, \infty\}$ is a cluster point of $A$,
and $b \in \R \cup \{ -\infty, \infty\}$ is a cluster point of $B$.
Suppose 
\begin{equation*}
\lim_{x \to a} f(x) = b\qquad \text{and} \qquad \lim_{y \to b} g(y) = c
\end{equation*}
for some $c \in \R \cup \{ -\infty, \infty \}$.
If $b \in B$, then suppose $g(b) = c$.
Then
\begin{equation*}
\lim_{x \to a} g\bigl(f(x)\bigr) = c .
\end{equation*}
\end{prop}

The proof is straightforward, and left as an exercise.  We already
know the proposition when $a, b, c \in \R$ (see Exercises
\ref{exercise:contlimitcomposition} and
\ref{exercise:contlimitbadcomposition}).  Again, the requirement that $g$ is
continuous at $b$, if $b \in B$, is necessary.

\begin{example}
Let $h(x) \coloneqq e^{-x^2+x}$.  Then
\begin{equation*}
\lim_{x\to \infty} h(x) = 0 .
\end{equation*}

Proof:
The claim follows once we know
\begin{equation*}
\lim_{x\to \infty} (-x^2+x) = -\infty
\end{equation*}
and
\begin{equation*}
\lim_{y\to -\infty} e^y = 0 ,
\end{equation*}
which is usually proved when the exponential function is defined.
\end{example}

\subsection{Exercises}

\begin{exercise}
Prove \propref{liminfty:unique}.
\end{exercise}

\begin{exercise}
Let $f \colon [1,\infty) \to \R$ be a function.  Define
$g \colon (0,1] \to \R$ via $g(x) \coloneqq f(\nicefrac{1}{x})$.
Using the definitions of limits directly,
show that $\lim_{x\to 0^+} g(x)$
exists if and only if $\lim_{x\to \infty} f(x)$ exists, in which
case they are equal.
\end{exercise}

\begin{exercise}
Prove \propref{prop:inflimcompositions}.
\end{exercise}

\begin{exercise}
Let us justify terminology.
Let $f \colon \R \to \R$ be a function such that
$\lim_{x \to \infty} f(x) = \infty$ (diverges to infinity).
Show that $f(x)$ diverges (i.e.\ does not converge) as $x \to \infty$.
\end{exercise}

\begin{exercise}
Come up with the definitions for limits of $f(x)$ going to $-\infty$ as $x \to
\infty$, $x \to -\infty$, and as $x \to c$ for a finite $c \in \R$.
Then state the definitions for limits of $f(x)$ going to $\infty$ 
as $x \to -\infty$, and as $x \to c$ for a finite $c \in \R$.
\end{exercise}

\begin{exercise}
Suppose $P(x) \coloneqq x^n + a_{n-1} x^{n-1} + \cdots + a_1 x + a_0$ is a \emph{\myindex{monic polynomial}}
of degree $n \geq 1$ (monic means that the coefficient of $x^n$ is 1).
\begin{enumerate}[a)]
\item
Show that if $n$ is even, then $\lim\limits_{x\to\infty} P(x) = 
\lim\limits_{x\to-\infty} P(x) = \infty$.
\item
Show that if $n$ is odd, then
$\lim\limits_{x\to\infty} P(x) = \infty$ and
$\lim\limits_{x\to-\infty} P(x) = -\infty$ (see previous exercise).
\end{enumerate}
\end{exercise}

\begin{exercise}
Let $\{ x_n \}_{n=1}^\infty$ be a sequence.  Consider $S \coloneqq \N \subset \R$, and
$f \colon S \to \R$ defined by $f(n) \coloneqq x_n$.  Show that
the two notions of limit,
\begin{equation*}
\lim_{n\to\infty} x_n \qquad \text{and} \qquad
\lim_{x\to\infty} f(x) 
\end{equation*}
are equivalent.  That is, show that if one exists so does
the other one, and in this case they are equal.
\end{exercise}

\begin{exercise}
Extend \lemmaref{seqflimitinf:lemma} as follows.
Suppose $S \subset \R$ has a cluster point $c \in \R$, $c = \infty$,
or $c = -\infty$.  Let $f \colon S \to \R$ be a function and suppose
$L = \infty$ or $L = -\infty$.  Show that
\begin{equation*}
\begin{aligned}
&
\lim_{x\to c} f(x) = L
\\
&
\qquad \text{if and only if}
\\
&
\lim_{n\to\infty} f(x_n) = L \enspace \text{for all sequences }
\{ x_n \}_{n=1}^\infty \text{ such that } \lim_{n\to\infty} x_n = c
\text{ and } x_n \in S \setminus \{c\}
\text{ for all } n  .
\end{aligned}
\end{equation*}
\end{exercise}

\begin{exercise}
Suppose $f \colon \R \to \R$ is a 2-periodic function, that is $f(x + 2) =
f(x)$ for all $x$.  Define $g \colon \R \to \R$ by 
$g(0) \coloneqq f(0)$ and for $x \neq 0$ by
\begin{equation*}
g(x) \coloneqq f\left(\frac{\sqrt{x^2+1}-1}{x}\right)
\end{equation*}
\begin{enumerate}[a)]
\item
Find the function $\varphi \colon (-1,1) \to \R$ such that
$g\bigl(\varphi(t)\bigr) = f(t)$, that is, $\varphi^{-1}(x) = 
\frac{\sqrt{x^2+1}-1}{x}$ for $x \neq 0$.
\item
Show that $f$ is continuous if and only if $g$ is continuous and
\begin{equation*}
\lim_{x \to \infty} g(x) = 
\lim_{x \to -\infty} g(x) = 
f(1) = f(-1) .
\end{equation*}
\end{enumerate}
\end{exercise}

%%%%%%%%%%%%%%%%%%%%%%%%%%%%%%%%%%%%%%%%%%%%%%%%%%%%%%%%%%%%%%%%%%%%%%%%%%%%%%

\sectionnewpage
\section{Monotone functions and continuity}
\label{sec:monotonefunc}

%mbxINTROSUBSECTION

\sectionnotes{1 lecture (optional, can safely be omitted unless
\sectionref{sec:ift} is also covered, requires \sectionref{sec:limitatinf})}

\begin{defn}
Let $S \subset \R$.
We say $f \colon S \to \R$ is \emph{\myindex{increasing}}
(resp.\  \emph{\myindex{strictly increasing}}) if $x,y \in S$ with
$x < y$ implies $f(x) \leq f(y)$ (resp.\ $f(x) < f(y)$).
We define
\emph{\myindex{decreasing}} and
\emph{\myindex{strictly decreasing}} in the same way by switching the
inequalities for $f$.

If a function is either increasing or decreasing, we say it is
\emph{monotone}\index{monotone function}.  If it is
strictly increasing or strictly decreasing, we say it is
\emph{strictly monotone}\index{strictly monotone function}.
\end{defn}

Sometimes \emph{\myindex{nondecreasing}}
(resp.\ \emph{\myindex{nonincreasing}}) is used
for increasing (resp.\ decreasing) function to emphasize it is not
strictly increasing (resp.\ strictly decreasing).

If $f$ is increasing, then $-f$ is decreasing and vice versa.  Therefore,
many results about monotone functions can just be proved for, say, increasing
functions, and the results follow easily for decreasing functions.

\subsection{Continuity of monotone functions}

One-sided limits for monotone functions are computed
by computing infima and suprema.

\begin{prop} \label{prop:monotlimits}
Let $S \subset \R$, $c \in \R$,
$f \colon S \to \R$ be increasing,
and
$g \colon S \to \R$ be decreasing.
If $c$ is a cluster point of $S \cap (-\infty,c)$, then
\begin{equation*}
\lim_{x \to c^-} f(x) = \sup \{ f(x) : x < c, x \in S \}
%mbxSTARTIGNORE
\qquad \text{and} \qquad
%mbxENDIGNORE
%mbxlatex \quad \text{and} \quad
\lim_{x \to c^-} g(x) = \inf \{ g(x) : x < c, x \in S \} .
\end{equation*}
If $c$ is a cluster point of $S \cap (c,\infty)$, then
\begin{equation*}
\lim_{x \to c^+} f(x) = \inf \{ f(x) : x > c, x \in S \}
%mbxSTARTIGNORE
\qquad \text{and} \qquad
%mbxENDIGNORE
%mbxlatex \quad \text{and} \quad
\lim_{x \to c^+} g(x) = \sup \{ g(x) : x > c, x \in S \} .
\end{equation*}
If $\infty$ is a cluster point of $S$, then
\begin{equation*}
\lim_{x \to \infty} f(x) = \sup \{ f(x) : x \in S \}
%mbxSTARTIGNORE
\qquad \text{and} \qquad
%mbxENDIGNORE
%mbxlatex \quad \text{and} \quad
\lim_{x \to \infty} g(x) = \inf \{ g(x) : x \in S \} .
\end{equation*}
If $-\infty$ is a cluster point of $S$, then
\begin{equation*}
\lim_{x \to -\infty} f(x) = \inf \{ f(x) : x \in S \}
%mbxSTARTIGNORE
\qquad \text{and} \qquad
%mbxENDIGNORE
%mbxlatex \quad \text{and} \quad
\lim_{x \to -\infty} g(x) = \sup \{ g(x) : x \in S \} .
\end{equation*}
\end{prop}

Namely, all the one-sided limits exist whenever they make
sense.  Therefore, for monotone functions, when we say
the left-hand limit $x \to c^-$
exists, we mean that $c$ is a cluster point of $S \cap (-\infty,c)$,
and same for the right-hand limit.

\begin{proof}
Let us assume $f$ is increasing, and we will show the first
equality.  The rest of the proof is very similar and is left as an
exercise.

Let $a \coloneqq \sup \{ f(x) : x < c, x \in S \}$.  If $a = \infty$,
then given an $M \in \R$, there exists an $x_M \in S$, $x_M < c$, such that $f(x_M) > M$. 
As $f$ is increasing, $f(x) \geq f(x_M) >  M$ for all $x \in S$ with $x > x_M$.
Take $\delta \coloneqq c-x_M > 0$ to satisfy the definition of the limit
being infinity.

Next suppose $a < \infty$.
Let $\epsilon > 0$ be given.  Because $a$ is the supremum and
$S \cap (-\infty,c)$ is nonempty, $a \in \R$ and
there exists an
$x_\epsilon \in S$,
$x_\epsilon < c$,
such that $f(x_\epsilon) > a-\epsilon$.  As $f$ is increasing,
if $x \in S$ and $x_\epsilon < x < c$, we have
$a-\epsilon < f(x_\epsilon) \leq f(x) \leq a$.  Let
$\delta \coloneqq c-x_\epsilon$.  Then for $x \in S \cap (-\infty,c)$
with $\sabs{x-c} < \delta$,
we have $\babs{f(x)-a} < \epsilon$.
\end{proof}

Suppose $f \colon S \to \R$ is increasing, $c \in S$, and
that both one-sided limits exist.
Since $f(x) \leq f(c) \leq f(y)$
whenever $x < c < y$, taking the limits we obtain
\begin{equation*}
\lim_{x \to c^-} f(x) \leq f(c) \leq \lim_{x \to c^+} f(x) .
\end{equation*}
Then $f$ is continuous at $c$ if and only if both limits are equal
to each other (and hence equal to $f(c)$).  See also
\propref{prop:onesidedlimits}.
See \figureref{fig:figinccont} to get an idea of what a discontinuity
looks like.


\begin{cor} \label{cor:continterval}
If $I \subset \R$ is an interval and $f \colon I \to \R$ is 
monotone and not constant, then $f(I)$ is an interval if and only if $f$
is continuous.
\end{cor}

Assuming $f$ is not constant is to avoid the technicality
that $f(I)$ is a single point: $f(I)$ is a single
point if and only if $f$ is constant.  A constant function is 
continuous.

\begin{proof}
Without loss of generality, suppose $f$ is increasing.

First suppose $f$ is continuous.  Take two points
$f(x_1), f(x_2)$ in $f(I)$ and suppose
$f(x_1) < f(x_2)$.
As $f$ is increasing, $x_1 < x_2$.  By the
\hyperref[IVT:thm]{intermediate value theorem},
given $y$ with $f(x_1) < y < f(x_2)$, we find
a $c \in (x_1,x_2) \subset I$ such that $f(c) = y$, so $y \in f(I)$. 
Hence, $f(I)$ is an interval.

Let us prove the reverse direction by contrapositive.
Suppose $f$ is not continuous at $c \in I$,
and that $c$ is not an endpoint of $I$.
Let
\begin{equation*}
%mbxlatex \begin{aligned}
a
%mbxlatex &
\coloneqq \lim_{x \to c^-} f(x) = \sup \bigl\{ f(x) : x \in I, x < c \bigr\} ,
%mbxSTARTIGNORE
\qquad
%mbxENDIGNORE
%mbxlatex \\
b
%mbxlatex &
\coloneqq \lim_{x \to c^+} f(x) = \inf \bigl\{ f(x) : x \in I, x > c \bigr\} .
%mbxlatex \end{aligned}
\end{equation*}
As $c$ is a discontinuity, $a < b$.
If $x < c$, then $f(x) \leq a$, and
if $x > c$, then $f(x) \geq b$.  Therefore,
no point
in $(a,b) \setminus \bigl\{ f(c) \bigr\}$ is in $f(I)$.
There exists $x_1 \in I$ with $x_1 < c$, so
$f(x_1) \leq a$, and there exists $x_2 \in I$ with $x_2 > c$,
so $f(x_2) \geq b$.  Both $f(x_1)$ and $f(x_2)$ are in $f(I)$,
but there are points in between them that are not in $f(I)$.
So $f(I)$ is not an interval.  See \figureref{fig:figinccont}.

When $c \in I$ is an endpoint, the proof is similar and is left as an exercise.
\end{proof}

\begin{myfigureht}
\myincludepdft{figinccont}{%
A diagram of a graph of a function with a jump discontinuity and the effect
on the image of the function.
The function is defined on an interval marked I.  Points x sub 1, c, and
x sub 2 are marked inside I and come in that order from left to right.
The function increases continuously until x equals c,
where the function approaches a, that is, the limit of f of x as x approaches
c from below is a.  The value of f at c is f of c, which is larger than a
and smaller than b, where b is the limit of f of x as x approaches c from
above.  Then the function keeps rising continuously until the end of the
interval I.  The image of f is split in three pieces.
The lowest piece is
an interval whose upper bound is a, and this interval contains
f of x sub 1.
The next piece is simply the point f of c.
The third, the highest, piece is an interval starting at b
and this interval contains f of x sub 2.}
\caption{Increasing function $f \colon I \to \R$ discontinuity at
$c$.\label{fig:figinccont}}
\end{myfigureht}

A striking property of monotone functions is that they cannot have
too many discontinuities.

\begin{cor} \label{cor:monotcountcont}
Let $I \subset \R$ be an interval and
$f \colon I \to \R$ be monotone.  Then $f$ has at most
countably many discontinuities.
\end{cor}

\begin{proof}
Let $E \subset I$ be the set of all discontinuities
that are not endpoints of $I$.  As there are
only two endpoints, it is enough to show that $E$ is countable.
Without loss of generality, suppose $f$ is increasing.
We will define an injection $h \colon E \to \Q$.
For each $c \in E$,
both one-sided limits of $f$ exist as $c$ is not an endpoint.
Let
\begin{equation*}
%mbxlatex \begin{aligned}
a
%mbxlatex &
\coloneqq \lim_{x \to c^-} f(x) = \sup \bigl\{ f(x) : x \in I, x < c \bigr\} ,
%mbxSTARTIGNORE
\qquad
%mbxENDIGNORE
%mbxlatex \\
b
%mbxlatex &
\coloneqq \lim_{x \to c^+} f(x) = \inf \bigl\{ f(x) : x \in I, x > c \bigr\} .
%mbxlatex \end{aligned}
\end{equation*}
As $c$ is a discontinuity, $a < b$.  
There exists a rational number $q \in (a,b)$, so let $h(c) \coloneqq q$.
Suppose $d \in E$ is another discontinuity.
If $d > c$, there
exists an $x \in I$ with $c < x < d$, and so $\lim_{x \to d^-} f(x) \geq b$.
Hence, the rational number we choose for $h(d)$ is different from $q$,
since $q=h(c) < b$ and $h(d) > b$.
Similarly if $d < c$.  After making such a choice for
every element of $E$, we have a 
one-to-one (injective) function into $\Q$.  Therefore, $E$ is countable.
\end{proof}

\begin{example} \label{example:countdiscont}
Denote the largest integer less than or equal to $x$
by $\lfloor x \rfloor$.
Define $f \colon [0,1] \to \R$ by
\begin{equation*}
f(x) \coloneqq
x +
\sum_{n=0}^{\lfloor 1/(1-x) \rfloor}
2^{-n} ,
\end{equation*}
for $x < 1$ and $f(1) \coloneqq 3$.
It is an exercise to show that $f$ is strictly increasing, bounded, and
has a discontinuity at all points $1-\nicefrac{1}{k}$ for $k \in \N$.  In particular,
there are countably many discontinuities, but the function is bounded and
defined on a closed bounded interval.  See \figureref{fig:countdiscont}.
\begin{myfigureht}
\myincludegraphics{increasing-discont-fig}{%
A graph of a function on the interval from 0 to 1
that is composed of straight line segments sloped upwards.  There is
a first segment going across half the graph starting at 
y equal to 1.5 and ending at y equal to 2.  Then there is an upward
jump discontinuity and another much shorter line segment
also sloped upwards.  This set of smaller and smaller upward jumps and
shorter and shorter upward sloped line segments keeps going until
at x equal to 1 we get to y equal to 3.}
\caption{Strictly increasing function on $[0,1]$ with countably many
discontinuities.\label{fig:countdiscont}}
\end{myfigureht}

Similarly, one can find an example of a monotone function discontinuous on a dense set
such as the rational numbers.  See the exercises.
\end{example}


\subsection{Continuity of inverse functions}


A strictly monotone function $f$ is one-to-one (injective).
To see this fact,
notice that if $x \neq y$, then we can assume $x < y$.  Either $f(x) <
f(y)$ if $f$ is strictly increasing or $f(x) > f(y)$ if $f$ is strictly
decreasing, so $f(x) \neq f(y)$.
Hence, $f$ must have an inverse $f^{-1}$ defined on its range.

\begin{prop} \label{prop:invcont}
If $I \subset \R$ is an interval and $f \colon I \to \R$ is strictly
monotone, then the inverse $f^{-1} \colon f(I) \to I$ is continuous.
\end{prop}

\begin{proof}
Suppose $f$ is strictly increasing.  The proof is almost
identical for a strictly decreasing function.
Since $f$ is strictly increasing, so is $f^{-1}$.  That is, if $f(x) <
f(y)$, then we must have $x < y$ and therefore
$f^{-1}\bigl(f(x)\bigr) < f^{-1}\bigl(f(y)\bigr)$.

Take $c \in f(I)$.
If $c$ is not a cluster point of $f(I)$, then $f^{-1}$ is continuous at $c$
automatically.  So let $c$ be a cluster point of $f(I)$.
Suppose both of the following one-sided limits exist, that is, they both
make sense:
\begin{align*}
x_0 & \coloneqq \lim_{y \to c^-} f^{-1}(y) =
\sup \bigl\{ f^{-1}(y) : y < c, y \in f(I) \bigr\}
=
\sup \bigl\{ x \in I : f(x) < c \bigr\} , \\
x_1 & \coloneqq \lim_{y \to c^+} f^{-1}(y) =
\inf \bigl\{ f^{-1}(y) : y > c, y \in f(I) \bigr\}
=
\inf \bigl\{ x \in I : f(x) > c \bigr\} .
\end{align*}
We have $x_0 \leq x_1$ as $f^{-1}$ is increasing.
For all $x \in I$ where $x > x_0$, we have $f(x) \geq c$.  As $f$ is strictly increasing,
we must have $f(x) > c$ for all $x \in I$ where $x > x_0$.  Therefore,
\begin{equation*}
\{ x \in I : x > x_0 \} \subset \bigl\{ x \in I : f(x) > c \bigr\}.
\end{equation*}
The infimum of the left-hand set is $x_0$, and the infimum of the right-hand
set is $x_1$, so we obtain $x_0 \geq x_1$.
So $x_1 = x_0$, and $f^{-1}$ is continuous at $c$.

If one of the one-sided limits does not exist (that is, when it does not
make sense to take the limit from that side), the argument is similar
and is left as an exercise.
\end{proof}

\begin{example}
The proposition does not require $f$ itself to be continuous.  Let
$f \colon \R \to \R$ be defined by
\begin{equation*}
f(x) \coloneqq
\begin{cases}
x & \text{if } x < 0, \\
x+1 & \text{if } x \geq 0. \\
\end{cases}
\end{equation*}
The function $f$ is not continuous at $0$.
The image of $I = \R$ is the set 
$(-\infty,0)\cup [1,\infty)$, not an interval.
Then $f^{-1} \colon (-\infty,0)\cup [1,\infty)
\to \R$ can be written as
\begin{equation*}
f^{-1}(y) =
\begin{cases}
y & \text{if } y < 0, \\
y-1 & \text{if } y \geq 1. 
\end{cases}
\end{equation*}
It is not difficult to see that $f^{-1}$ is a continuous function.  See
\figureref{invcontfig} for the graphs.
\begin{myfigureht}
\myincludepdft{invcontfigAB}{%
Two graphs.  On the left, an upward sloping diagonal line that goes from
bottom left to the origin (not including the origin),
then jumps up and then continues diagonally up.
On the right graph, the first line giving the
graph is the same, but then the graph skips a short
interval and the line starts again on the horizontal
axis going diagonally up again.}
\caption{Graph of $f$ on the left and $f^{-1}$ on the right.\label{invcontfig}}
\end{myfigureht}
\end{example}

Notice what happens with the proposition if $f(I)$ is an interval.
In that case, we could simply
apply \corref{cor:continterval} to both $f$ and $f^{-1}$.  That is, if
$f \colon I \to J$ is an onto strictly monotone function and $I$ and $J$ are intervals,
then both $f$ and $f^{-1}$ are continuous.  Furthermore, $f(I)$ is an
interval precisely when $f$ is continuous.

\subsection{Exercises}

\begin{exercise}
Suppose $f \colon [0,1] \to \R$ is monotone.  Prove $f$ is bounded.
\end{exercise}

\begin{exercise}
Finish the proof of \propref{prop:monotlimits}.
Hint: You can halve your work by noticing that if $g$ is decreasing,
then $-g$ is increasing.
\end{exercise}

\begin{exercise}
Finish the proof of \corref{cor:continterval}.
\end{exercise}

\begin{exercise}
Prove the claims in \exampleref{example:countdiscont}.
\end{exercise}

\begin{exercise}
Finish the proof of \propref{prop:invcont}.
\end{exercise}

\begin{samepage}
\begin{exercise}
Suppose $S \subset \R$, and $f \colon S \to \R$ is an increasing
function.  Prove:
\begin{enumerate}[a)]
\item
If $c$ is a cluster point
of $S \cap (c,\infty)$, then
$\lim\limits_{x\to c^+} f(x) < \infty$.
\item
If $c$ is a cluster point of $S \cap (-\infty,c)$
and $\lim\limits_{x\to c^-} f(x) = \infty$, then
$S \subset (-\infty,c)$.
\end{enumerate}
\end{exercise}
\end{samepage}

\begin{exercise}
Let $I \subset \R$ be an interval and $f \colon I \to \R$ a function.
Suppose that for each $c \in I$, there exist $a, b \in \R$ with
$a > 0$ such that $f(x) \geq a x + b$ for all $x \in I$
and $f(c) = a c + b$.  Show that $f$ is strictly increasing.
\end{exercise}

\begin{exercise}
Suppose $I$ and $J$ are intervals and
$f \colon I \to J$ is a continuous, bijective (one-to-one and onto)
function.  Show that $f$ is strictly monotone.
\end{exercise}

\begin{exercise}
Consider a monotone function $f \colon I \to \R$ on an interval $I$.  Prove that there exists
a function $g \colon I \to \R$ such that
$\lim\limits_{x \to c^-} g(x) = g(c)$ for all $c$ in $I$ except the
smaller (left) endpoint of $I$, and such that
$g(x) = f(x)$ for all but countably many $x \in I$.
\end{exercise}

\begin{exercise}
\leavevmode
\begin{enumerate}[a)]
\item
Let $S \subset \R$ be a subset.  If $f \colon S \to \R$ is increasing and
bounded,
then show that there exists an increasing $F \colon \R \to \R$
such that $f(x) = F(x)$ for all $x \in S$.
\item
Find an example of a strictly increasing bounded $f \colon S \to \R$ such that
an increasing $F$ as above cannot be strictly increasing no matter how we
choose it.
\end{enumerate}
\end{exercise}

\begin{exercise}[Challenging] \label{exercise:increasingfuncdiscatQ}
Find an example of an increasing function $f \colon [0,1] \to \R$
that has a discontinuity at each rational number.  Then show that the image
$f\bigl([0,1]\bigr)$ contains no interval.  Hint: Enumerate
the rational numbers and define
the function with a series.
\end{exercise}

\begin{exercise}
Suppose $I$ is an interval and $f \colon I \to \R$ is monotone.
Show that $\R \setminus f(I)$ is a countable union of disjoint intervals.
\end{exercise}

\begin{exercise}
Suppose $f \colon [0,1] \to (0,1)$ is increasing.  Show that for every
$\epsilon > 0$, there exists
a strictly increasing $g \colon [0,1] \to (0,1)$ such that
$g(0) = f(0)$, $f(x) \leq g(x)$ for all $x$, and $g(1)-f(1) < \epsilon$.
\end{exercise}

\begin{exercise}
Prove that the \myindex{Dirichlet function}
$f \colon [0,1] \to\R$, defined by $f(x) \coloneqq 1$
if $x$ is rational and $f(x) \coloneqq 0$ otherwise, cannot be written as a
difference of two increasing functions.  That is, there do not exist
increasing $g$ and $h$ such that $f(x) = g(x) - h(x)$.
\end{exercise}

\begin{exercise}
Suppose $f \colon (a,b) \to (c,d)$ is a strictly increasing
onto function.  Prove that there exists a $g \colon (a,b) \to (c,d)$,
which is also strictly increasing and onto, and $g(x) < f(x)$ for all $x \in
(a,b)$.
\end{exercise}


